\documentclass[sigconf]{cidr-2027}

\usepackage{microtype}
\usepackage{enumitem}
\usepackage{xspace}
\usepackage{booktabs}
\usepackage{tikz}
\usepackage{colortbl}

\setlist{nosep,leftmargin=1.25em}

\newtheorem{theorem}{Theorem}
\newtheorem{proposition}{Proposition}
\newtheorem{lemma}{Lemma}
\newtheorem{corollary}{Corollary}
\newtheorem{definition}{Definition}
\theoremstyle{remark}
\newtheorem{example}{Example}
\newtheorem{remark}{Remark}

% Use \paragraph as minor section-like header
\makeatletter
\renewcommand\paragraph{\@startsection{paragraph}{4}{\z@}%
  {-.5\baselineskip \@plus -2\p@ \@minus -.2\p@}%
  {-0.5em}%
  {\normalfont\bfseries}}
\makeatother

\newcommand{\InputOrd}{\preceq_I}
\newcommand{\OutputOrd}{\preceq_O}
\newcommand{\InputGe}{\succeq_I}
\newcommand{\OutputGe}{\succeq_O}
% Output event emitting #2 after applying the \commitment sequence #1
% (\varepsilon = no \commitment).
\newcommand{\Out}[2]{\mathsf{output}_{#1}(#2)}
\newcommand{\commitment}{commitment\xspace}
\newcommand{\commitments}{commitments\xspace}
\newcommand{\Commitment}{Commitment\xspace}
\newcommand{\Commitments}{Commitments\xspace}
% Order vocabulary (role-based names; change the words here).
% Input order:  "j \inrel i" means i \InputOrd j (non-strict); inputs \ingrow;
%   a common \indesc \inrel each input (articles: "an \indesc").
% Output order: "u \outrel o" means o \OutputOrd u; a common \outnoun
%   \outrel each output.
\newcommand{\inrel}{extends\xspace}          % relation, 3rd person
\newcommand{\inrelbase}{extend\xspace}       % relation, base form
\newcommand{\inrelpart}{extending\xspace}    % relation, participle
\newcommand{\inrelpast}{extended\xspace}     % relation, past participle
\newcommand{\ingrows}{grows\xspace}          % growth, 3rd person
\newcommand{\ingrow}{grow\xspace}            % growth, base form
\newcommand{\inmore}{more extended\xspace}   % informal comparative
\newcommand{\indesc}{input extension\xspace} % noun
\newcommand{\indescs}{input extensions\xspace}
\newcommand{\Indesc}{Input extension\xspace}
\newcommand{\outrel}{extends\xspace}          % relation, 3rd person
\newcommand{\outrelbase}{extend\xspace}       % relation, base form
\newcommand{\outrelpart}{extending\xspace}    % relation, participle
\newcommand{\outrelpast}{extended\xspace}     % relation, past participle
\newcommand{\outnoun}{output extension\xspace}% noun
\newcommand{\outnouns}{output extensions\xspace}
\setlength{\emergencystretch}{1.5em}
\newcommand{\Spec}{\mathcal{S}}
% The name of the specification model; optional suffix supports \spectype[s].
\newcommand{\spectype}[1][]{I/O specification#1\xspace}
\newcommand{\Spectype}[1][]{I/O specification#1\xspace}
\newcommand{\fiber}[1]{P(#1)}
\newcommand{\dog}{\texttt{dog}}
\newcommand{\catv}{\texttt{cat}}
\newcommand{\Subsumed}{\mathit{Subsumed}}
\newcommand{\Excluded}{\mathit{Excluded}}
\newcommand{\grayrow}[1]{\textcolor{gray}{#1}}
\newcommand{\sh}[1]{\cellcolor{gray!22}#1}
% Full-width strike line through the row just ended; place after \\.
\newcommand{\strikerow}{%
  \noalign{\vskip-0.5\dimexpr\ht\strutbox+\dp\strutbox\relax}%
  \noalign{\vskip-0.5\arrayrulewidth}%
  \arrayrulecolor{black!65}\hline\arrayrulecolor{black}%
  \noalign{\vskip 0.5\dimexpr\ht\strutbox+\dp\strutbox\relax}%
  \noalign{\vskip-0.5\arrayrulewidth}}
% Pointer to the long version (proofs and extended material).
\newcommand{\extended}{the full paper~\cite{hellerstein2026completecalmfull}}
\newcommand{\fullproof}{Proof in the full paper~\cite{hellerstein2026completecalmfull}.}
\newcommand{\proofinfull}{\begin{proof}In the full paper~\cite{hellerstein2026completecalmfull}.\end{proof}}

% ---------- Comment macro; NEVER EVER DELETE!!! ----------
\newcommand{\jmh}[1]{\textcolor{red}{\textit{#1 -- JMH}}}


\usepackage{xcolor}
\makeatletter
\newcommand{\cxtag}[2]{\textsuperscript{\textsf{\tiny[#1\,#2]}}}
% \cxdef{id}{s|e}{x sp}{y sp}{page}: read back from the .aux file.
\newcommand{\cxdef}[5]{\expandafter\gdef\csname cx@#1@#2\endcsname{{#3}{#4}{#5}}}
\newcommand{\cx@unpack}[4]{\def#1##1##2##3{%
  \dimen@=##1sp\edef#2{\strip@pt\dimen@}\dimen@=##2sp\edef#3{\strip@pt\dimen@}%
  \def#4{##3}}}
\newcommand{\cx@geom}[1]{% #1: page number
  \ifodd#1\relax\pgfmathsetmacro\cx@L{72.27+\oddsidemargin}%
  \else\pgfmathsetmacro\cx@L{72.27+\evensidemargin}\fi
  \pgfmathsetmacro\cx@T{\paperheight-72.27-\topmargin-\headheight-\headsep}%
  \pgfmathsetmacro\cx@B{\cx@T-\textheight}}
\newcommand{\cx@col}[2]{% x, macro to set
  \pgfmathsetmacro#2{((#1-\cx@L)>(\columnwidth+0.5*\columnsep))?1:0}}
\def\cx@rgb#1 #2 #3\relax{\pgfsys@color@rgb@stroke{#1}{#2}{#3}}
\newcommand{\cx@line}[4]{% "r g b", x, y1, y2 (pt, page coordinates)
  \begin{tikzpicture}[remember picture,overlay]
    \cx@rgb#1\relax
    \draw[line width=2.5pt] ([shift={(#2pt,#3pt)}]current page.south west)
      -- ([shift={(#2pt,#4pt)}]current page.south west);
  \end{tikzpicture}}
% \cxbar{id}{color}{s|e}: save this position; draw the bar piece on this page.
\newcommand{\cxbar}[3]{%
  \pdfsavepos
  \write\@auxout{\string\cxdef{#1}{#3}{\the\pdflastxpos}{\the\pdflastypos}{\thepage}}%
  \ifcsname cx@#1@s\endcsname\ifcsname cx@#1@e\endcsname
    \begingroup
    \expandafter\expandafter\expandafter\cx@gets\csname cx@#1@s\endcsname
    \expandafter\expandafter\expandafter\cx@gete\csname cx@#1@e\endcsname
    \cx@geom{\ifx#3s\cx@sp\else\cx@ep\fi}%
    \cx@col{\cx@sx}\cx@cs \cx@col{\cx@ex}\cx@ce
    \pgfmathsetmacro\cx@same{(\cx@sp==\cx@ep && \cx@cs==\cx@ce && \cx@ey<\cx@sy)?1:0}%
    \pgfmathsetmacro\cx@top{\cx@sy+8}\pgfmathsetmacro\cx@bot{\cx@ey-2}%
    \if#3s%
      \pgfmathsetmacro\cx@X{\cx@L+\cx@cs*(\columnwidth+\columnsep)-6}%
      \ifdim\cx@same pt=1pt \cx@line{#2}{\cx@X}{\cx@top}{\cx@bot}%
      \else \cx@line{#2}{\cx@X}{\cx@top}{\cx@B}\fi
    \else
      \pgfmathsetmacro\cx@X{\cx@L+\cx@ce*(\columnwidth+\columnsep)-6}%
      \ifdim\cx@same pt=1pt \else \cx@line{#2}{\cx@X}{\cx@T}{\cx@bot}\fi
    \fi
    \endgroup
  \fi\fi}
\cx@unpack\cx@gets\cx@sx\cx@sy\cx@sp
\cx@unpack\cx@gete\cx@ex\cx@ey\cx@ep
\makeatother
\begin{document}

\title{Complete CALM: A Unified Criterion for Coordination}

\author{Joseph M. Hellerstein}
\affiliation{%
  \institution{UC Berkeley \& Amazon Web Services}
  \city{Berkeley, CA}
  \country{USA}
}

\color{orange!90!black}\begin{abstract}
\cxtag{1}{O/L}When does a distributed system intrinsically need coordination?  The literature
has lacked a common formalism for the question.  CALM, CAP, and CRDTs each answer it
for one class of programs, consistency models, or design patterns.  Their
answers cannot be compared or combined, and together they leave basic questions
open, such as how coordination repairs inconsistency.  We put them on common
ground.  We define a minimalist \spectype: a
possibility table that includes, for each input, all outputs a system may
produce.  Each column carries its own partial order, saying when one input
\inrel another and when one output \outrel another.  An \spectype is
\emph{consistent} when all its possible runs are confluent: runs whose inputs have a common
\indesc can always catch up to a common state.  Possibilities narrow as input
grows or outputs are emitted. An irrevocable \emph{\commitment} removes possibilities not subsumed by those
remaining.  Coordination mechanisms such as locks and consensus enforce these
\commitments.

\color{orange!90!black}\cxtag{2}{O/L}In this unified setting, \emph{Functional CALM} generalizes
the CALM Theorem from Datalog queries to any function between ordered domains:
such a specification is consistent exactly when it is monotone: no
\commitment is required. We then go further: \emph{Complete CALM} handles nondeterministic specifications
that allow several outputs for the same input, where the test becomes
\emph{Lifted Monotonicity}: at every input, some permitted output \outrel
whatever could have been emitted up to that point.
\emph{Complete CAP} is a corollary: the CAP trade-off arises
exactly when Lifted Monotonicity fails.
\emph{Proper Coordination} identifies which stage of a composition must make
the \commitments, and certifies that later stages need none of their own.  Highly Available Transactions, Invariant
Confluence, and later refinements of CALM all map into the same framework.
\end{abstract}

\begin{CCSXML}
<ccs2012>
<concept>
<concept_id>10003752.10003790.10003795</concept_id>
<concept_desc>Theory of computation~Distributed computing models</concept_desc>
<concept_significance>500</concept_significance>
</concept>
<concept>
<concept_id>10002951.10003317.10003347.10003356</concept_id>
<concept_desc>Information systems~Database transaction processing</concept_desc>
<concept_significance>300</concept_significance>
</concept>
</ccs2012>
\end{CCSXML}

\ccsdesc[500]{Theory of computation~Distributed computing models}
\ccsdesc[300]{Information systems~Database transaction processing}
\keywords{coordination, monotonicity, CALM, CAP, transactions, CRDTs}

\maketitle
\noindent{\small\textsf{Text color, benefit of a rewrite: \textcolor{red!80!black}{R large}, \textcolor{orange!90!black}{O moderate}, \textcolor{blue!75!black}{B minor polish}, N leave (black).  Margin bar, intrinsic complexity: \textcolor{red!80!black}{XL}, \textcolor{orange!90!black}{L}, \textcolor{blue!60}{M}, S no bar.  Tags: [paragraph rewrite/complexity].}}\par

% ======================================================================
\section{Introduction}
\label{sec:intro}

\color{black}\cxtag{3}{N/S}Concurrent systems rely on coordination mechanisms ranging from mutexes and
atomic instructions to database locks and consensus protocols. These mechanisms
prevent concurrent actions from producing undesirable outcomes.  They also bring
costs: overhead, latency, and availability in distributed settings.  Avoiding
unnecessary coordination is therefore a central goal in distributed systems and
databases.  But when is coordination avoidable, and when is it intrinsically
required? Various answers are scattered across the literature in incompatible
forms. We unify those results via a common formalism, and use it to address new
questions: when do nondeterministic specifications require coordination, what does
coordination achieve, and how do we know when it is properly applied?

\color{black}\cxtag{4}{N/M}\cxbar{4}{0.4 0.4 1}{s}The CALM Conjecture proposed a crisp, bidirectional criterion:
consistency can be achieved without coordination exactly when a specification
is monotone~\cite{hellerstein2010declarative}.  But the original proof applied
only in a restricted setting: Datalog queries evaluated by relational transducer
networks~\cite{ameloot2013relational,hellerstein2020keeping}. A related
question comes from the CAP conjecture: which consistency guarantees can remain
available when participants cannot communicate? There, the original proof
applied in a different restricted setting: linearizable
registers~\cite{brewer2000towards,gilbert2002cap}.  Both results are invoked
quite broadly: CALM as a general rule that monotone programs need no
coordination, CAP as a general constraint on distributed systems. The folklore
outgrew the proofs.\cxbar{4}{0.4 0.4 1}{e}

\color{blue!75!black}\cxtag{5}{B/M}\cxbar{5}{0.4 0.4 1}{s}The apparent connection between CALM and CAP was never established in a common
formalism, and nearby results only complicate the picture. CRDTs use
inflationary updates over join-semilattices for coordination-free
convergence~\cite{shapiro2011crdt}; HATs identify transactional guarantees
available under partition~\cite{bailis2013hat}; invariant confluence asks when
independently reached database states merge
safely~\cite{bailis2014coordination}. The diverse formalisms leave cross-cutting
questions unclear, including whether CALM encompasses CAP, whether CRDT join
laws are necessary, why only some transactional isolation levels remain
available during partition, and whether invariant confluence differs from CALM
(\S\ref{sec:instances}).\cxbar{5}{0.4 0.4 1}{e}

\color{black}\cxtag{6}{N/S}We now turn to unifying these ideas under a single formalism.

\paragraph{The Model: \Spectype[s].}
\color{black}\cxtag{7}{N/S}In systems, as in life, our possibilities narrow over time. Sometimes the world
moves on, and possibilities are lost to us. Sometimes possibilities disappear
because we must decide among them. A decision is incidental when
every possibility it removes is covered by an alternative that remains. A
decision is fateful when it rules out a possibility that no remaining
alternative can cover. When must a system make a fateful decision?

\color{blue!75!black}\cxtag{8}{B/M}\cxbar{8}{0.4 0.4 1}{s}We ground the problem in a simple model of \emph{\spectype[s]}, or \emph{specs}
for short, independent of any language, protocol, or placement. The designer of
a spec provides a \emph{possibility relation} $P\subseteq I\times O$: a virtual
table that may be infinite, like a mathematical function. Each row in this table
is a possibility observable from the outside: it pairs the input received so far
with a permitted output. The design includes the input domain $I$ and output domain $O$,
with partial orders that help capture the spec's semantics.  On inputs, the order
is accumulation: $i\InputOrd j$ means that $j$ \emph{\inrel} $i$ with more input
from the environment.  On outputs, the order is subsumption: $o\OutputOrd u$
means that  $u$ \emph{\outrel} $o$ by subsuming all spec-relevant information in
$o$. These are both application-specific notions that must be captured well for
a spec to be faithful.\cxbar{8}{0.4 0.4 1}{e}

\color{black}\cxtag{9}{N/M}\cxbar{9}{0.4 0.4 1}{s}State-machine-style specs can be projected into this form, skipping
operational details that are not externally observable (\S\ref{sec:scope}).
Despite its simplicity, our model captures not only application invariants,
but also general-purpose consistency definitions such as linearizability,
causal consistency, and transactional isolation levels
(\S\ref{sec:complete-cap}, \S\ref{sec:hat-instance}, and \extended).\cxbar{9}{0.4 0.4 1}{e}

\color{black}\cxtag{10}{N/M}\cxbar{10}{0.4 0.4 1}{s}As a system runs, possibilities narrow: the spec's table loses rows.
Newly-arrived input removes rows that cannot be reached by extending that input;
the world has moved on.  Coordination mechanisms such as locks and consensus
remove rows by decision.  We call their relevant semantic effect a
\emph{\commitment}: an irrevocable, fateful system decision.  It can remove a permitted
output that no remaining output at the same input covers: for example, choosing between two competing candidates
in a leader election protocol removes the possibility of the loser winning that
election.  Or it can rule out all rows with a given input; for example, closing
the polls for an election rules out consideration of any future ballots.\cxbar{10}{0.4 0.4 1}{e}

\paragraph{The Criterion: Lifted Monotonicity.}
\color{blue!75!black}\cxtag{11}{B/M}\cxbar{11}{0.4 0.4 1}{s}With this model in hand, we are ready to define the key aspects of CALM:
Consistency and Monotonicity. A spec is \emph{consistent} when
all its possible executions are \emph{confluent}: states reached along
different paths can be brought to a common state. We will see that this definition
is general, accommodating narrower notions from prior literature.
Our monotonicity criterion requires more care: it must capture the full
expressivity of our model. It comes via two invariants. The first is always
required; the second addresses nondeterministic specs.\cxbar{11}{0.4 0.4 1}{e}

\color{orange!90!black}\cxtag{12}{O/L}\cxbar{12}{0.9 0.45 0}{s}\emph{Learn and refine.} As execution progresses, input grows: the system learns more from
the environment.
For specs that permit only one output at each input, the possibility table is a function
and monotonicity is standard: $i\InputOrd j$ implies $f(i)\OutputOrd f(j)$.
\emph{Functional CALM} (\S\ref{sec:functional-calm-theorem}) proves that consistency is equivalent to
monotonicity in this setting.\cxbar{12}{0.9 0.45 0}{e}
% Emitting a monotone function's output is always an incidental choice:
% the output at any \inrelpast input is \outrelpart: it covers the earlier output.
% This setting generalizes Datalog-based CALM papers as well as CRDT state specifications.

% If monotonicity
% fails, some output requires a fateful decision: a \commitment that rules
% out later input which could contradict it.  The result covers the relational query
% languages of prior CALM papers, and the usual CRDT state specification.

\color{blue!75!black}\cxtag{13}{B/M}\cxbar{13}{0.4 0.4 1}{s}\emph{Choose and align.} Prior work on CALM and CRDTs addresses functional
settings. \emph{Nondeterministic specs}, including leader election and transaction
scheduling, are not functions: many outputs are within spec for a given input.
This requires an additional monotonicity invariant.  When many runs of a system
choose different outputs for the same input, those choices
must \outrelbase together---the outputs are \emph{directed}, in order theory terms.
Together, monotonicity under input growth and output choice make all output decisions incidental.
We call the pair of invariants \emph{Lifted Monotonicity}.
\emph{Complete CALM} proves that consistency is equivalent to Lifted
Monotonicity (\S\ref{sec:complete-calm-theorem}).\cxbar{13}{0.4 0.4 1}{e}

\color{black}\cxtag{14}{N/M}\cxbar{14}{0.4 0.4 1}{s}What of coordination?  Locks, agreement protocols, and other coordination
mechanisms enforce \commitments.  Our criterion identifies the \commitments a
spec's outputs require, independently of the mechanism used.
Distributed histories are captured under the same consistency criteria as centralized ones
(\S\ref{sec:distributed-execution}). Complete CAP (\S\ref{sec:complete-cap}) says
Lifted Monotonicity exactly characterizes specs that can keep every
permitted output available and confluent during a partition.  The
Gilbert--Lynch theorem follows directly.\cxbar{14}{0.4 0.4 1}{e}

\paragraph{Contributions.}
\color{black}\cxtag{15}{N/S}We make contributions in three categories:

\color{blue!75!black}\noindent
\cxtag{16}{B/M}\cxbar{16}{0.4 0.4 1}{s}\emph{1. A unified model.} A system's observable behavior is captured via
        \spectype[s]. Consistency is defined as
        confluence of the runs a spec allows, unifying general consistency models,
        application-specific requirements, and their distributed executions.\cxbar{16}{0.4 0.4 1}{e}

\color{black}\noindent
\cxtag{17}{N/M}\cxbar{17}{0.4 0.4 1}{s}\emph{2. A unified criterion.} For functions between ordered domains,
        Functional CALM shows that consistency is equivalent to monotonicity. For nondeterministic specs,
        we prove Complete CALM: Consistency As Lifted Monotonicity. Complete CAP is a corollary:
        the CAP tradeoff holds only under failure of Lifted Monotonicity.\cxbar{17}{0.4 0.4 1}{e}

\color{blue!75!black}\noindent
\cxtag{18}{B/M}\cxbar{18}{0.4 0.4 1}{s}\emph{3. \Commitments, Composition and Proper Coordination.} We define \commitments to capture how
        systems exclude possibilities that are not covered by alternatives. We then show how specs
        compose by equijoin, so that, in composition, upstream
        stages can free downstream stages from issuing \commitments.\cxbar{18}{0.4 0.4 1}{e}

\color{black}\vspace{0.5em}
\noindent
\cxtag{19}{N/S}The full paper~\cite{hellerstein2026completecalmfull} gives further
intuition, proofs, and applications.

% ======================================================================
\section{Functional CALM}
\label{sec:functional-calm}

\color{black}\cxtag{20}{N/S}This section focuses on functional specs, in which each input has a
single specified output.  We walk through the framework informally on an
election (\S\ref{sec:election-story}), state it precisely
(\S\ref{sec:framework-brief}), and
prove Functional CALM (\S\ref{sec:functional-calm-theorem}), identifying
when an output can be emitted without \commitment.

\subsection{The framework by example}
\label{sec:election-story}

\begin{figure*}[t]
  \centering
  \small
  \newcommand{\rowhdr}{\toprule $\mathrm{in}$ & $\mathrm{out}$ \\ \midrule}
  \setlength{\tabcolsep}{6pt}
  \begin{tabular}[t]{@{}l c@{}}
    \multicolumn{2}{@{}l}{(a) Initial state, current input $a\cdot\cdot$}\\
    \rowhdr
    \sh{$\mathbf{a\cdot\cdot}$} & \sh{$\bot$} \\
    $ab\cdot$ & \sh{$\bot$} \\
    $aba$ & \sh{$a$} \\
    $abb$ & \sh{$b$} \\
    $aa\cdot$ & \sh{$a$} \\
    $aaa$ & \sh{$a$} \\
    $aab$ & \sh{$a$} \\
    $a\cdot a$ & \sh{$a$} \\
    $a\cdot b$ & \sh{$\bot$} \\
    \bottomrule
    \multicolumn{2}{@{}l}{\footnotesize future-output set $\{\bot,a,b\}$; undetermined}\\
    \multicolumn{2}{@{}l}{\footnotesize $\bot$ needs no \commitment}
  \end{tabular}
  \hfill
  \begin{tabular}[t]{@{}l c@{}}
    \multicolumn{2}{@{}l}{(b) After $\mathsf{input}(ab\cdot)$}\\
    \rowhdr
    \grayrow{$a\cdot\cdot$} & \grayrow{$\bot$} \\
    \sh{$\mathbf{ab\cdot}$} & \sh{$\bot$} \\
    $aba$ & \sh{$a$} \\
    $abb$ & \sh{$b$} \\
    \grayrow{$aa\cdot$} & \grayrow{$a$} \\
    \grayrow{$aaa$} & \grayrow{$a$} \\
    \grayrow{$aab$} & \grayrow{$a$} \\
    \grayrow{$a\cdot a$} & \grayrow{$a$} \\
    \grayrow{$a\cdot b$} & \grayrow{$\bot$} \\
    \bottomrule
    \multicolumn{2}{@{}l}{\footnotesize future-output set $\{\bot,a,b\}$; undetermined}\\
    \multicolumn{2}{@{}l}{\footnotesize $\bot$ needs no \commitment}
  \end{tabular}
  \hfill
  \begin{tabular}[t]{@{}l c@{}}
    \multicolumn{2}{@{}l}{(c) After $\mathsf{input}(aa\cdot)$ instead}\\
    \rowhdr
    \grayrow{$a\cdot\cdot$} & \grayrow{$\bot$} \\
    \grayrow{$ab\cdot$} & \grayrow{$\bot$} \\
    \grayrow{$aba$} & \grayrow{$a$} \\
    \grayrow{$abb$} & \grayrow{$b$} \\
    \sh{$\mathbf{aa\cdot}$} & \sh{$a$} \\
    $aaa$ & \sh{$a$} \\
    $aab$ & \sh{$a$} \\
    \grayrow{$a\cdot a$} & \grayrow{$a$} \\
    \grayrow{$a\cdot b$} & \grayrow{$\bot$} \\
    \bottomrule
    \multicolumn{2}{@{}l}{\footnotesize future-output set $\{a\}$; determined}\\
    \multicolumn{2}{@{}l}{\footnotesize $a$ needs no \commitment}
  \end{tabular}
  \hfill
  \begin{tabular}[t]{@{}l c@{}}
    \multicolumn{2}{@{}l}{(d) (b), then polls closed}\\
    \rowhdr
    \grayrow{$a\cdot\cdot$} & \grayrow{$\bot$} \\
    \sh{$\mathbf{ab\cdot}$} & \sh{$\bot$} \\
    \grayrow{$aba$} & \grayrow{$a$} \\ \strikerow
    \grayrow{$abb$} & \grayrow{$b$} \\ \strikerow
    \grayrow{$aa\cdot$} & \grayrow{$a$} \\
    \grayrow{$aaa$} & \grayrow{$a$} \\
    \grayrow{$aab$} & \grayrow{$a$} \\
    \grayrow{$a\cdot a$} & \grayrow{$a$} \\
    \grayrow{$a\cdot b$} & \grayrow{$\bot$} \\
    \bottomrule
    \multicolumn{2}{@{}l}{\footnotesize future-output set $\{\bot\}$; determined}\\
    \multicolumn{2}{@{}l}{\footnotesize $\bot$ is the result}
  \end{tabular}
  \caption{The election specification and three states reached from it, limited to rows
  whose input \inrel $a\cdot\cdot$.  Bold: current input; shaded: the current
  input and its future-output set.  Gray: lost after an input event; struck
  through if removed by a \commitment.  Panels (b)
  and (c) follow input events.  In (d) the registrar closes the polls, a \commitment that emitting
  $\bot$ does not require, and emits $\bot$.}
  \label{fig:election-frames}
\end{figure*}

\color{black}\cxtag{21}{N/S}We introduce the framework with a small election, illustrated in
Figure~\ref{fig:election-frames}.  Each row of the possibility table pairs
the ballots received so far with the permitted result; the column orders
say how inputs \ingrow and which later outputs \outrelbase earlier ones.

\color{black}\cxtag{22}{N/S}Three voters $v_1,v_2,v_3$ each cast at most one immutable ballot for $a$ or
$b$.  The output is whichever candidate has two ballots, or $\bot$ while there
is no winner.  An input is the ballots received so far, written as a
three-character string with a dot for a ballot not yet received.  For brevity,
the figure shows only inputs \inrelpart $a\cdot\cdot$.

\color{blue!75!black}\cxtag{23}{B/M}\cxbar{23}{0.4 0.4 1}{s}The two orders require thoughtful design.  The \emph{input order}
must record accumulated ballots, one per voter: $ab\cdot$ \inrel $a\cdot\cdot$, while
$aa\cdot$ and $ab\cdot$ are incomparable.  The \emph{output order} says which
announcements are safe to build on: $\{\bot \OutputOrd a, \bot \OutputOrd b\}$.
$\bot$ claims nothing, so either winner
\outrel it, while $a$ and $b$ are incomparable.  In general, the functional table is \emph{monotone}
when
\[
  i\InputOrd j\quad\Longrightarrow\quad f(i)\OutputOrd f(j).
\]
Here, for example, the input sequence $a\cdot\cdot\InputOrd ab\cdot\InputOrd aba$
monotonically produces the output sequence $\bot\OutputOrd\bot\OutputOrd a$.\cxbar{23}{0.4 0.4 1}{e}

\paragraph{Input and determination.}
\color{orange!90!black}\cxtag{24}{O/L}\cxbar{24}{0.9 0.45 0}{s}Each panel of Figure~\ref{fig:election-frames} captures the specification state: a pair of a \emph{current input}
and a \emph{remainder table} of rows
still possible from that input.  When $v_2$'s ballot for $b$ arrives, the input event
$\mathsf{input}(ab\cdot)$ produces panel~(b), causing the table to passively lose rows whose inputs no
longer \inrelbase the ballots received; such inputs are now illegal.  The
\emph{future-output set} consists of the outputs in the remaining rows at the
current input and its \indescs: in panel~(b),
$\{\bot,a,b\}$.
If $v_2$ instead votes for $a$ (panel~(c)), the future-output set is $\{a\}$.  The state
is \emph{determined}; in this case input alone resolved the outcome, without \commitment.\cxbar{24}{0.9 0.45 0}{e}

\color{blue!75!black}\cxtag{25}{B/L}\cxbar{25}{0.9 0.45 0}{s}Panel~(d), by contrast, involves the \emph{system}: the implementation that
receives the ballots and announces the result
(\S\ref{sec:systems}).  In this example, the system acts as a registrar: when it emits
$\bot$ at $ab\cdot$, it also closes the polls.  Its irrevocable \emph{\commitment}
to close the polls removes the still-possible $aba$ and $abb$ rows but leaves the current input
unchanged, and so makes $\bot$ the final result.  Emitting $\bot$ by itself
would need no \commitment, since every output still possible \outrel it;
closing the polls is a further fateful decision by the system, not something the
spec requires.\cxbar{25}{0.9 0.45 0}{e}

\begin{figure}[t]
  \centering
  \small
  \setlength{\tabcolsep}{6pt}
  \begin{minipage}{0.86\columnwidth}
  (e) As (b), all outputs incomparable, $\bot$ emitted\par\smallskip
  \begin{minipage}[t]{0.44\linewidth}
    \begin{tabular*}{\linewidth}[t]{@{}l@{\extracolsep{\fill}}c@{}}
      \toprule
      $\mathrm{in}$ & $\mathrm{out}$ \\ \midrule
      \grayrow{$a\cdot\cdot$} & \grayrow{$\bot$} \\
      \sh{$\mathbf{ab\cdot}$} & \sh{$\bot$} \\
      \grayrow{$aba$} & \grayrow{$a$} \\ \strikerow
      \grayrow{$abb$} & \grayrow{$b$} \\ \strikerow
      \grayrow{$aa\cdot$} & \grayrow{$a$} \\
    \end{tabular*}
  \end{minipage}\hfill
  \begin{minipage}[t]{0.44\linewidth}
    \begin{tabular*}{\linewidth}[t]{@{}l@{\extracolsep{\fill}}c@{}}
      \toprule
      $\mathrm{in}$ (cont.) & $\mathrm{out}$ (cont.) \\ \midrule
      \grayrow{$aaa$} & \grayrow{$a$} \\
      \grayrow{$aab$} & \grayrow{$a$} \\
      \grayrow{$a\cdot a$} & \grayrow{$a$} \\
      \grayrow{$a\cdot b$} & \grayrow{$\bot$} \\
      \bottomrule
    \end{tabular*}
    {\footnotesize future-output set $\{\bot\}$; determined}
  \end{minipage}\par\smallskip
  \end{minipage}
  \caption{Panel (b) under the equality order, after $\mathsf{output}(\bot)$, wrapped
  to preserve space.  The registrar's \commitment in
  Figure~\ref{fig:election-frames}(d) is required here.}
  \label{fig:nonmonotone-frame}
\end{figure}

\paragraph{Output and consistency.}
\color{black}\cxtag{26}{N/M}\cxbar{26}{0.4 0.4 1}{s}In a functional table, an output is a \emph{legal output} when every output in
the future-output set \outrel it.  Hence $\bot$ is legal
in panels~(a) and~(b), and $a$ in panel~(c).
The output event $\mathsf{output}(o)$ emits $o$ and keeps only the rows whose
outputs extend $o$, since later outputs must not contradict it.  When $o$ is
legal, every remaining row qualifies and nothing is lost.  When it is not, the
event removes the rows that conflict with $o$, ruling out their inputs; that
removal is a \commitment.\cxbar{26}{0.4 0.4 1}{e}

\color{blue!75!black}\cxtag{27}{B/M}\cxbar{27}{0.4 0.4 1}{s}Figure~\ref{fig:nonmonotone-frame} considers a variant of our example that
changes only the output order: all outputs are incomparable, including $\bot$
for no winner (the \emph{equality order}, under which each element is related only
to itself).  Then $\bot$ is incompatible with future winners, so emitting it
requires a \commitment.  That change in the output order is Functional CALM in
miniature: the original order is monotone, so later outputs \outrelbase every
earlier one; under the equality order, some outputs can be emitted
only with a \commitment.\cxbar{27}{0.4 0.4 1}{e}

\color{blue!75!black}\cxtag{28}{B/L}\cxbar{28}{0.9 0.45 0}{s}The need to issue a \commitment is a property solely of the table and its orders.  Rows $(ab\cdot,\bot)$ and
$(aba,a)$ have the common \indesc $aba$, but under the equality
order no output at $aba$ \outrel both outputs $\bot$ and $a$.  A run that emits
$\bot$ at $ab\cdot$ and one that receives $aba$ cannot be reconciled at $aba$.\cxbar{28}{0.9 0.45 0}{e}

\color{black}\cxtag{29}{N/S}Input and output events form a transition system, and the spec is
consistent when its runs are confluent.
Section~\ref{sec:functional-calm-theorem} shows that for functional tables
this corresponds exactly to monotonicity.

\subsection{The functional framework in brief}
\label{sec:framework-brief}
\label{sec:functional-specifications}
\label{sec:functional-events}
\label{sec:executions}

\color{black}\cxtag{30}{N/S}We now make the functional framework precise: functional I/O
specs and monotonicity, then a transition system that models their
runs, and finally consistency as confluence of those runs.

\paragraph{Specs and monotonicity.} \color{blue!75!black}\cxtag{31}{B/M}\cxbar{31}{0.4 0.4 1}{s}We begin with specs.
\begin{definition}[Functional \spectype, total]
\label{def:functional-specification}
A \emph{functional \spectype} is a tuple
$\Spec=(I,\InputOrd,O,\OutputOrd,P)$,
where $(I,\InputOrd)$ and $(O,\OutputOrd)$ are partially ordered input and
output types, and $P\subseteq I\times O$ is its \emph{possibility relation}:
a typed relation with schema $(\mathrm{in},\mathrm{out})$ in which every
input appears in at least one row ($P$ is \emph{total} on $I$) and at most one
(a functional dependency).  The unique output paired with $i$ induces a
function $f:I\to O$.\cxbar{31}{0.4 0.4 1}{e}
\end{definition}

\color{orange!90!black}\cxtag{32}{O/L}\cxbar{32}{0.9 0.45 0}{s}Because an output may depend on everything a system has received, we model
inputs as \emph{cumulative}: $i\InputOrd j$ means that $j$ contains at least
the input information in $i$.  We say that $j$ \emph{\inrel} $i$, or is an
\emph{\indesc} of $i$.  We call $\OutputOrd$ the \emph{output order}:
$o\OutputOrd u$ means that $u$ subsumes every externally visible claim in $o$
and may add more, and we say that $u$ \emph{\outrel} $o$.  A \emph{common
\indesc} $j$ of $i_1,\ldots,i_k$ is an input that \inrel every $i_r$, and a
\emph{common \outnoun} of $o_1,\ldots,o_k$ is an output that \outrel each $o_r$.  Neither need be
unique or least: for example, two independent transactions $u$ and $v$, with schedules
ordered by subsequence, have two incomparable minimal common \indescs, $uv$
and $vu$.\cxbar{32}{0.9 0.45 0}{e}

\color{black}\begin{definition}[Functional monotonicity]
\label{def:functional-monotonicity}
\cxtag{33}{N/S}$\Spec$ is \emph{monotone} if
\begin{equation}
  i\InputOrd j
  \quad\Longrightarrow\quad
  f(i)\OutputOrd f(j).
  \label{eq:functional-monotonicity}
\end{equation}
\end{definition}

\color{black}\cxtag{34}{N/M}\cxbar{34}{0.4 0.4 1}{s}We introduce some compact notation. For a relation $R\subseteq I\times O$, a predicate in square brackets selects
rows, and commas separate conjuncts: $R[\mathrm{in}\InputGe j,\
\mathrm{out}\OutputGe o]$ holds the rows whose input \inrel $j$ and whose
output \outrel $o$.  We say that a row $(k,u)$ is \emph{at} $k$, and $\mathrm{out}(R)$ is
the set of outputs in $R$.\cxbar{34}{0.4 0.4 1}{e}

\color{black}\cxtag{35}{N/S}We next define the transition system.

\paragraph{States: remaining possibilities.}
\color{blue!75!black}\cxtag{36}{B/L}\cxbar{36}{0.9 0.45 0}{s}As illustrated by the non-gray text in Figure~\ref{fig:election-frames}, a
\emph{specification state} $\sigma=(i,R)$ records the current input and the
possibilities still open: $i$ is the \emph{current input}, and $R\subseteq P$
is the \emph{remainder table}.  The table $R$ contains a row at $i$ and only
rows at \indescs of $i$.  An input $k$ is \emph{legal} when $R$ contains a
row at $k$.  Rows are never added during execution.  The \emph{canonical state
at $i$} has received $i$, emitted nothing, and removed no rows:
\begin{equation}
  \sigma_P(i)=\bigl(i,\ P[\mathrm{in}\InputGe i]\bigr).
  \label{eq:canonical-state}
\end{equation}
The \emph{fiber} of $R$ at $k$ is the set of outputs paired with $k$:
$R(k)=\mathrm{out}(R[\mathrm{in}=k])$.  The \emph{future-output set} of
$(i,R)$ is $\mathrm{out}(R)$: all outputs still possible from $i$.  The state
is \emph{determined} when its future-output set has one element.\cxbar{36}{0.9 0.45 0}{e}

\paragraph{Events: narrowing possibilities.}
\color{black}\cxtag{37}{N/M}\cxbar{37}{0.4 0.4 1}{s}Each event applies a selection predicate on one column of the remainder table.
For a state $(i,R)$ and a legal input $j$, the \emph{input event}
\begin{equation}
  (i,R)\xrightarrow{\mathsf{input}(j)}
  \bigl(j,\ R[\mathrm{in}\InputGe j]\bigr)
  \label{eq:input-event}
\end{equation}
makes $j$ current and drops the rows whose input does not
\inrelbase $j$.  Because inputs are cumulative, $\mathsf{input}(j)$ then
$\mathsf{input}(k)$ has the effect of $\mathsf{input}(k)$ alone.\cxbar{37}{0.4 0.4 1}{e}

\color{black}\cxtag{38}{N/S}A system controls what it tells the world.  Later outputs must \outrelbase
earlier ones.

\color{blue!75!black}\begin{definition}[Legal output]
\label{def:exposable}
\cxtag{39}{B/M}\cxbar{39}{0.4 0.4 1}{s}At a state $(i,R)$, let $o$ be the unique output at the current input.  The
output $o$ is \emph{legal} when every output in the future-output set \outrel
$o$.\cxbar{39}{0.4 0.4 1}{e}
\end{definition}

\color{black}\cxtag{40}{N/S}Selecting on the output column may require a \commitment.

\color{orange!90!black}\begin{definition}[Output event, canonical \commitment]
\label{def:output-event}
\cxtag{41}{O/L}\cxbar{41}{0.9 0.45 0}{s}At a state $(i,R)$ with current output $o$, the \emph{output event}
$\mathsf{output}(o)$ atomically keeps the rows whose outputs \outrelbase $o$,
and emits $o$:
\begin{equation}
  (i,R)\xrightarrow{\mathsf{output}(o)}\bigl(i,\ R[\mathrm{out}\OutputGe o]\bigr).
  \label{eq:output-action}
\end{equation}
If $o$ is not legal, the transition applies the \emph{canonical \commitment}
$x_o$.  For every row it removes, no row remains at the same input.\cxbar{41}{0.9 0.45 0}{e}
\end{definition}

\color{orange!90!black}\cxtag{42}{O/L}\cxbar{42}{0.9 0.45 0}{s}Thus $o$ is legal exactly when $\mathsf{output}(o)$ applies no \commitment.  A
\commitment removes rows that input alone would preserve; afterward, their
inputs are no longer legal.  Systems may issue more \commitments than an output
requires, as the registrar does in Figure~\ref{fig:election-frames}(d), but
such extra \commitments are a system's choice, not the spec's
(\S\ref{sec:systems}).\cxbar{42}{0.9 0.45 0}{e}

\color{black}\cxtag{43}{N/S}An event from a determined state always leads to another determined state.

\paragraph{Runs and consistency.}
\color{black}\cxtag{44}{N/M}\cxbar{44}{0.4 0.4 1}{s}We now consider the executions corresponding to a spec.
\begin{definition}[Specification event, run, \commitment-free]
\label{def:run}
Input and output events are \emph{specification events}, or just \emph{events}.  A \emph{run} is a
finite sequence of them from the
canonical state $\sigma_P(i)$ at some input $i$.  It is
\emph{\commitment-free} when no output event in it applies a \commitment; the
states such runs reach are \emph{reachable without \commitment}.\cxbar{44}{0.4 0.4 1}{e}
\end{definition}

\color{black}\cxtag{45}{N/S}We compare a system's choices with every behavior the declared spec permits,
not only the behaviors that system happens to produce.

\color{blue!75!black}\cxtag{46}{B/M}\cxbar{46}{0.4 0.4 1}{s}Input events keep only rows whose input \inrel the current one, and output
events keep only rows whose outputs \outrelbase what has been emitted.  Both
steps are relational selections, and selections commute.\cxbar{46}{0.4 0.4 1}{e}

\color{black}\begin{lemma}[Rows remaining after a run]
\label{lem:run-residual}
\cxtag{47}{N/M}\cxbar{47}{0.4 0.4 1}{s}If a run ends with the current input $i$ after emitting
$o_1,\ldots,o_m$, its remainder table is \(P[\mathrm{in}\InputGe i,\ \mathrm{out}\OutputGe o_1,\ \ldots,\
  \mathrm{out}\OutputGe o_m].\)\cxbar{47}{0.4 0.4 1}{e}
\end{lemma}
\proofinfull

\color{black}\cxtag{48}{N/M}\cxbar{48}{0.4 0.4 1}{s}Different runs may retain different rows.  At a common \indesc $j$ of their
current inputs, they can meet only if they still share a possible row at $j$.\cxbar{48}{0.4 0.4 1}{e}

\color{blue!75!black}\begin{definition}[Consistency, confluence at an input]
\label{def:consistency}
\cxtag{49}{B/L}\cxbar{49}{0.9 0.45 0}{s}An \spectype is \emph{consistent} when every finite family of
runs is confluent: if the runs end at states
$(i_1,R_1),\ldots,(i_k,R_k)$ and $j$ is a common \indesc of their current
inputs, each can be continued to a
common state with current input $j$.  They are then \emph{confluent at} $j$.\cxbar{49}{0.9 0.45 0}{e}
\end{definition}

\color{blue!75!black}\cxtag{50}{B/L}\cxbar{50}{0.9 0.45 0}{s}Consistency requires confluence at every common \indesc, not merely at some
later one.  Otherwise conflicting outputs can escape before reconciliation.
In the equality-order election, $f(ab\cdot)=\bot$ and $f(aba)=a$ are
incomparable.  One run can emit $\bot$ at $ab\cdot$, while another accepts
$aba$ and can emit $a$.  The runs cannot reconcile at their common \indesc
$aba$, and waiting cannot retract either output from the environment.  More
generally, even when later input could reconcile the internal states, the
environment may already have acted on their outputs through channels the system
does not see, as in Lamport's ``anomalous behavior''~\cite{lamport1978time};
the reconciling input may never arrive.\cxbar{50}{0.9 0.45 0}{e}

\color{black}\cxtag{51}{N/S}In fact, sharing such a row is also sufficient.

\color{black}\begin{proposition}[Survivor test]
\label{prop:confluence}
\cxtag{52}{N/L}\cxbar{52}{0.9 0.45 0}{s}Runs ending at states $(i_r,R_r)$, $r=1,\ldots,k$, are confluent
at $j$ if and only if some row $(j,u)$ at $j$ is in every $R_r$.  For any such
$u$, $\mathsf{input}(j)$ followed by $\mathsf{output}(u)$ takes every final
state to $\bigl(j,\ P[\mathrm{in}\InputGe j,\ \mathrm{out}\OutputGe u]\bigr)$.\cxbar{52}{0.9 0.45 0}{e}
\end{proposition}

\color{black}\begin{proof}[Proof sketch]
\cxtag{53}{N/M}\cxbar{53}{0.4 0.4 1}{s}Rows are never added, so the row of a common state at $j$ was there
in every run.  \fullproof\cxbar{53}{0.4 0.4 1}{e}
\end{proof}

\color{blue!75!black}\cxtag{54}{B/M}\cxbar{54}{0.4 0.4 1}{s}Consistency ensures that runs never diverge irreparably: no ``split brain''.  Whatever input
and output events a finite collection of runs has taken, any input \inrelpart all
of theirs leaves a single state that every run can reach and that contradicts no
prior output.
Consistency is a property of the spec.  It describes no protocol and
promises no delivery or termination.\cxbar{54}{0.4 0.4 1}{e}

\subsection{The Functional CALM theorem}
\label{sec:functional-calm-theorem}

\color{black}\cxtag{55}{N/S}Consistency is defined by a spec's executions, monotonicity by its
table and orders alone.  Functional CALM shows that these coincide.  Full proofs are in \extended.

\color{black}\begin{theorem}[Functional CALM]
\label{thm:functional-calm}
\cxtag{56}{N/S}A functional \spectype is consistent if and only if it is
monotone.
\end{theorem}

\color{blue!75!black}\begin{proof}[Proof sketch]
\cxtag{57}{B/L}\cxbar{57}{0.9 0.45 0}{s}A monotone $f$ only emits outputs that all of its future outputs \outrelbase,
so the row $(j,f(j))$ at a common \indesc $j$ survives in every run.
If $f$ is not monotone, a run that emits $f(i)$ and a run that accepts an
\indesc $k$ of $i$ with $f(i)\not\OutputOrd f(k)$ can never be reconciled at $k$.
\fullproof\cxbar{57}{0.9 0.45 0}{e}
\end{proof}

\color{black}\begin{corollary}[State is fixed by input]
\label{cor:state-by-input}
\cxtag{58}{N/M}\cxbar{58}{0.4 0.4 1}{s}In a consistent functional specification, every run ending at current
input $i$ is in the canonical state $\sigma_P(i)$, whatever outputs it emitted
along the way.
\end{corollary}
\begin{proof}[Proof sketch]
By monotonicity, the outputs a run emits remove no row at an input extension of
$i$.
\fullproof\cxbar{58}{0.4 0.4 1}{e}
\end{proof}

\color{black}\cxtag{59}{N/S}Knowing the input is therefore enough to know the state, and with it the one
output.  Monotonicity also governs legality:

\color{black}\begin{proposition}[Output-legality boundary]
\label{prop:functional-exposure-boundary}
\cxtag{60}{N/M}\cxbar{60}{0.4 0.4 1}{s}A functional \spectype is monotone if and only if every output at
every state reachable without \commitment is legal.\cxbar{60}{0.4 0.4 1}{e}
\end{proposition}

\color{blue!75!black}\begin{proof}[Proof sketch]
\cxtag{61}{B/M}\cxbar{61}{0.4 0.4 1}{s}By Corollary~\ref{cor:state-by-input}, monotone $f$ reaches only canonical
states. At $\sigma_P(i)$, $f(i)$ is legal exactly
when $f$ is monotone at $i$.
 \fullproof\cxbar{61}{0.4 0.4 1}{e}
\end{proof}

\color{blue!75!black}\cxtag{62}{B/M}\cxbar{62}{0.4 0.4 1}{s}A monotone spec can thus emit every current output without \commitment. If a
spec is non-monotone, then for some $i\InputOrd k$,
$f(i)\not\OutputOrd f(k)$, so at $\sigma_P(i)$ the event
$\mathsf{output}(f(i))$ must exclude $k$.\cxbar{62}{0.4 0.4 1}{e}

\color{blue!75!black}\cxtag{63}{B/M}\cxbar{63}{0.4 0.4 1}{s}This is the familiar CALM dichotomy stated at the level of the spec
rather than inside a particular programming language.
Classical CALM fits, with inputs $J$ as sets of facts ordered by
inclusion and output views $Q(J)$ ordered by inclusion.  Functional
monotonicity becomes $J\subseteq J'\Rightarrow Q(J)\subseteq Q(J')$,
the standard definition of a monotone query~\cite{abiteboul1995foundations};
positive Datalog satisfies it, and negation may not.  Functional CALM
generalizes the statement to arbitrary input and
output orders, whatever formalism expresses the spec.  CRDT state specs have the same functional shape,
with received payloads ordered by inclusion and states by the semilattice
order~\cite{shapiro2011crdt}; Sections~\ref{sec:query-instance}
and~\ref{sec:crdt-instance} treat both.\cxbar{63}{0.4 0.4 1}{e}

% ======================================================================
\section{Complete CALM}
\label{sec:complete-calm}

\color{black}\cxtag{64}{N/S}We now extend CALM theory to nondeterministic specs, whose fibers may contain
several outputs.  Section~\ref{sec:lifted-monotonicity} defines Lifted
Monotonicity, Section~\ref{sec:lifted-output-events} lifts output events, and
Section~\ref{sec:complete-calm-theorem} proves Complete CALM.

\subsection{Nondeterministic specs and Lifted Monotonicity}
\label{sec:ambiguity-story}
\label{sec:relational-specifications}
\label{sec:lifted-monotonicity}

\color{blue!75!black}\cxtag{65}{B/M}\cxbar{65}{0.4 0.4 1}{s}We begin with our most general definition of a spec, which admits
output nondeterminism.
\begin{definition}[\Spectype, nondeterminism]
\label{def:relational-specification}
An \emph{\spectype} is
$\Spec=(I,\InputOrd,O,\OutputOrd,P)$ with $P\subseteq I\times O$,
where $(I,\InputOrd)$ and $(O,\OutputOrd)$ are partially ordered sets.  The
relation $P$ is the spec's \emph{possibility relation} and is total
on $I$.  A row $(i,o)\in P$ is a possible input/output pair: output $o$ is
permitted at input $i$.  The fiber
$\fiber{i}=\{o\in O:(i,o)\in P\}$ may now contain several
outputs; the \spectype is \emph{nondeterministic} when some fiber does.  A
functional \spectype is the special case $\fiber{i}=\{f(i)\}$.\cxbar{65}{0.4 0.4 1}{e}
\end{definition}

\color{blue!75!black}\cxtag{66}{B/M}\cxbar{66}{0.4 0.4 1}{s}Specification nondeterminism describes permitted results, not a strategy's method of choosing
among them, and does not by itself imply a \commitment. Specification states and input events
(\S\ref{sec:framework-brief}) apply unchanged to relational tables.\cxbar{66}{0.4 0.4 1}{e}

\color{black}\cxtag{67}{N/M}\cxbar{67}{0.4 0.4 1}{s}Figure~\ref{fig:ambiguity-frames} contrasts two specs with several outputs
in one fiber: serial schedules for two transactions, and route sets observed by a
traceroute-like tool.  For the discussion here, transaction outputs are simply
serial schedules, to which executions correspond via equivalence;
Section~\ref{sec:hat-instance} studies more thorough transactional specs.\cxbar{67}{0.4 0.4 1}{e}

\begin{figure}[t]
  \centering
  \small
  \setlength{\tabcolsep}{4pt}
  \begin{minipage}[t]{0.36\columnwidth}
    \raggedright
    (f) After emitting $[A,B]$\par\smallskip
    \begin{tabular*}{\linewidth}[t]{@{}l@{\extracolsep{\fill}}l@{}}
      \toprule
      $\mathrm{in}$ & $\mathrm{out}$ \\ \midrule
      $\{A,B\}$ & \sh{$[A,B]$} \\
      \grayrow{$\{A,B\}$} & \grayrow{$[B,A]$} \\ \strikerow
      \bottomrule
    \end{tabular*}\par\smallskip
    {\footnotesize no remaining row \outrel the removed one: a \commitment}
  \end{minipage}\hfill
  \begin{minipage}[t]{0.6\columnwidth}
    \raggedright
    (g) After emitting $\{s{\to}\texttt{s1}\}$\par\smallskip
    \begin{tabular*}{\linewidth}[t]{@{}l@{\extracolsep{\fill}}l@{}}
      \toprule
      $\mathrm{in}$ & $\mathrm{out}$ \\ \midrule
      $i$ & \sh{$\{s{\to}\texttt{s1}\}$} \\
      \grayrow{$i$} & \grayrow{$\{s{\to}h\}$} \\
      \grayrow{$i$} & \grayrow{$\{s{\to}h,\ s{\to}h{\to}\texttt{s2}\}$} \\
      $i$ & \sh{$\{s{\to}\texttt{s1},\ s{\to}h\}$} \\
      $i$ & \sh{$\{s{\to}\texttt{s1},\ s{\to}h,\ s{\to}h{\to}\texttt{s2}\}$} \\
      \bottomrule
    \end{tabular*}\par\smallskip
    {\footnotesize each removed row is \outrelpast by a remaining one: an output
    subsumption}
  \end{minipage}
  \caption{Two specs with several rows at one input, each after one output
  event.  (f) Serial orders of two concurrent transactions, under the equality
  order.  (g) Route sets from a source $s$ to a service on \texttt{s1} and
  \texttt{s2}, ordered by inclusion; $s$ reaches \texttt{s1} directly and
  \texttt{s2} through hop $h$.  Legend as in Figure~\ref{fig:election-frames};
  here gray alone marks output subsumption.}
  \label{fig:ambiguity-frames}
\end{figure}

\color{blue!75!black}\cxtag{68}{B/L}\cxbar{68}{0.9 0.45 0}{s}The transaction schedules $[A,B]$ and $[B,A]$ are incompatible, so emitting
either requires a \commitment that removes the other.  In the traceroute spec,
every removed route set is \outrelpast by a remaining one that also \outrelbase
the emitted observation.  Its removal is \emph{output subsumption}, not
\commitment.  Lifted Monotonicity combines two properties: Relational
Monotonicity requires each output to retain some \outnoun as input \ingrows,
while directedness requires the outputs in each fiber to have common \outnouns.
We define them in turn.\cxbar{68}{0.9 0.45 0}{e}

\color{black}\begin{definition}[Relational Monotonicity]
\label{def:relational-monotonicity}
\cxtag{69}{N/M}\cxbar{69}{0.4 0.4 1}{s}An \spectype satisfies \emph{Relational Monotonicity}
when every output keeps some \outnoun as input \ingrows:
        \cxbar{69}{0.4 0.4 1}{e}\begin{equation}
  i\InputOrd j,\ o\in\fiber{i}
  \quad\Longrightarrow\quad
  \exists u\in\fiber{j}:\ o\OutputOrd u.
  \label{eq:relational-monotonicity}
        \end{equation}
\end{definition}

\color{blue!75!black}\cxtag{70}{B/M}\cxbar{70}{0.4 0.4 1}{s}For singleton fibers this is functional monotonicity,
Equation~\ref{eq:functional-monotonicity}.  Relational Monotonicity treats each
output separately.  Directedness supplies the second requirement.  For a set
$F\subseteq O$ of outputs, outputs $o_1,\ldots,o_k$ are
\emph{compatible in $F$} if some $u\in F$ \outrel all of them, and
\emph{incompatible in $F$} otherwise.\cxbar{70}{0.4 0.4 1}{e}



\color{black}\begin{definition}[Directed fiber, directed fibers]
\label{def:fiber-directedness}
\cxtag{71}{N/M}\cxbar{71}{0.4 0.4 1}{s}A fiber $F\subseteq O$ is \emph{directed} when every nonempty finite
$X\subseteq F$ is compatible in $F$.  An \spectype has \emph{directed fibers} when every fiber
$\fiber{i}$ is directed.\cxbar{71}{0.4 0.4 1}{e}
\end{definition}

\color{blue!75!black}\cxtag{72}{B/L}\cxbar{72}{0.9 0.45 0}{s}A singleton fiber is trivially directed.  As with common \indescs, a
common \outnoun need not be
least: directedness requires some upper bound within the fiber for each finite
set, not a least one or a fixed merge operator, so it is weaker than a
join-semilattice.  For example, the fiber
\[
  \{\{a\},\{b\},\{a,b,c\},\{a,b,d\},\{a,b,c,d\}\}
\]
under inclusion is directed, since $\{a,b,c,d\}$ \outrel every member, yet $\{a\}$ and $\{b\}$ have two
incomparable minimal common \outnouns.  Thus
the traceroute fiber is directed: the union of finitely many route sets in
it is again in the fiber and extends each. The two transaction orders in
Figure~\ref{fig:ambiguity-frames}(f), by contrast, are incompatible.\cxbar{72}{0.9 0.45 0}{e}

\paragraph{Lifted Monotonicity.}
\color{black}\begin{definition}[Lifted Monotonicity, monotone specification]
\label{def:lifted-monotonicity}
\cxtag{73}{N/M}\cxbar{73}{0.4 0.4 1}{s}An \spectype satisfies \emph{Lifted Monotonicity}, and is called
\emph{monotone}, when, for every $i\InputOrd j$ and every nonempty finite
$X\subseteq\fiber{i}$, some $u\in\fiber{j}$ \outrel every member of $X$.\cxbar{73}{0.4 0.4 1}{e}
\end{definition}

\color{blue!75!black}\cxtag{74}{B/S}With singleton fibers it is functional monotonicity
(Definition~\ref{def:functional-monotonicity}). Lifted
Monotonicity holds exactly when Relational Monotonicity and directed fibers
both do.

\paragraph{The common-extension condition.}
\color{black}\cxtag{75}{N/M}\cxbar{75}{0.4 0.4 1}{s}Lifted Monotonicity compares an input with its
own input extensions. Runs may emit outputs at incomparable inputs, say
$o_1$ at $i_1$ and $o_2$ at $i_2$, and can be reconciled only at a common
\indesc $j$ of both. The next condition asks exactly that: a common
\outnoun of $o_1$ and $o_2$ in the fiber at $j$.\cxbar{75}{0.4 0.4 1}{e}

\color{black}\begin{definition}[Common-extension condition]
\label{def:common-expansion-condition}
\cxtag{76}{N/M}\cxbar{76}{0.4 0.4 1}{s}An \spectype satisfies the \emph{common-extension condition} when,
for every finite set of rows $(i_1,o_1),\ldots,(i_k,o_k)$ of $P$ and every
common \indesc $j$ of $i_1,\ldots,i_k$, some $u\in P(j)$ \outrel every
$o_r$.\cxbar{76}{0.4 0.4 1}{e}
\end{definition}

\color{black}\cxtag{77}{N/S}For the empty set of rows the condition asks only that $P(j)$ be nonempty, which
follows from $P$ being total on $I$.

\color{black}\begin{proposition}[Common-extension form of Lifted Monotonicity]
\label{prop:common-expansion-form}
\cxtag{78}{N/S}An \spectype satisfies Lifted Monotonicity if and only
if it satisfies the common-extension condition.
\end{proposition}

\color{blue!75!black}\begin{proof}[Proof sketch]
\cxtag{79}{B/M}\cxbar{79}{0.4 0.4 1}{s}Relational Monotonicity carries each $o_r$ to some
\outnoun in the fiber at $j$, and directedness of that fiber gives those a
common \outnoun there.  \fullproof\cxbar{79}{0.4 0.4 1}{e}
\end{proof}

\subsection{Lifted output events}
\label{sec:lifted-output-events}

\color{red!80!black}\cxtag{80}{R/XL}\cxbar{80}{0.8 0 0}{s}Lifted Monotonicity is a property of the possibility relation and its orders.  We now return to execution.
For a state $(i,R)$ and $o\in R(i)$, the output event $\mathsf{output}(o)$
again emits $o$ and keeps exactly the rows whose outputs \outrelbase $o$,
ending at $(i,R[\mathrm{out}\OutputGe o])$ as in Equation~\ref{eq:output-action}.
Every later output must \outrelbase $o$, so the event removes the rows whose
outputs do not.  Each removed row $(k,x)$ falls into one of two classes:
subsumed or excluded.  A row $(k,x)$ is placed in $\Subsumed_o$ if some
$u\in R(k)$ \outrel both $x$ and $o$.  Then a later output $u$ reports
everything $x$ would have, and the removal, an output subsumption, loses
nothing.  Otherwise $(k,x)$ is placed in $\Excluded_o$, and the canonical
\commitment $x_o$ removes it.  If $x_o$ removes every row at $k$, it rules out
$k$; otherwise it rules out only those outputs.  In the functional
case, $\Subsumed_o$ is empty.  Formal definitions are in \extended.\cxbar{80}{0.8 0 0}{e}

\color{blue!75!black}\cxtag{81}{B/L}\cxbar{81}{0.9 0.45 0}{s}Returning to Figure~\ref{fig:nonmonotone-frame},
emitting $\bot$ under the equality order
puts the rows for $aba$ and $abb$ in $\Excluded_o$ and rules out those inputs;
emitting $[A,B]$ in Figure~\ref{fig:ambiguity-frames}(f) puts the row for
$[B,A]$ in $\Excluded_o$ and rules out that output at $\{A,B\}$; and the
route sets removed in panel~(g) are in $\Subsumed_o$.\cxbar{81}{0.9 0.45 0}{e}

\color{blue!75!black}\cxtag{82}{B/L}\cxbar{82}{0.9 0.45 0}{s}With several outputs at one input, legality is more permissive: an output
event may remove rows by output subsumption without applying a \commitment.
As before, $o$ is a legal output when $\mathsf{output}(o)$ applies no
\commitment, exactly when $\Excluded_o$ is empty.
Equivalently, at every legal \indesc $j$ of $i$, $o$ and each
remaining output at $j$ have a common \outnoun at $j$.\cxbar{82}{0.9 0.45 0}{e}

\subsection{The Complete CALM theorem}
\label{sec:complete-calm-theorem}

\color{orange!90!black}\cxtag{83}{O/M}\cxbar{83}{0.4 0.4 1}{s}A lifted output event ends at $(i,R[\mathrm{out}\OutputGe o])$, as in
Equation~\ref{eq:output-action}.  States, input events, runs,
\commitment-free runs, and consistency are defined as before.  The results of
Section~\ref{sec:executions} depend on output events only through that equation,
so the rows-remaining lemma (Lemma~\ref{lem:run-residual}) and survivor test
(Proposition~\ref{prop:confluence}) apply unchanged.  Thus consistency reduces
to Lifted Monotonicity.\cxbar{83}{0.4 0.4 1}{e}

\color{black}\begin{theorem}[Complete CALM]
\label{thm:complete-calm}
\cxtag{84}{N/S}An \spectype is consistent if and only if it satisfies Lifted
Monotonicity, equivalently the common-extension condition.
\end{theorem}

\color{blue!75!black}\begin{proof}[Proof sketch]
\cxtag{85}{B/M}\cxbar{85}{0.4 0.4 1}{s}If the common-extension condition holds, a common \outnoun $u\in P(j)$
of everything emitted survives in every run.
If outputs emitted at inputs with a common \indesc $j$ are incompatible in
$P(j)$, the runs that emit them share no row at $j$ and can never be
reconciled there.  \fullproof\cxbar{85}{0.4 0.4 1}{e}
\end{proof}

\color{blue!75!black}\cxtag{86}{B/M}\cxbar{86}{0.4 0.4 1}{s}The name CALM, for Consistency As Logical Monotonicity, tied the idea to a
logic language.  The new, more general result is Consistency As Lifted Monotonicity.
The condition that characterizes consistency is now a
property of a spec: Lifted Monotonicity, which for a function is
ordinary monotonicity.  Section~\ref{sec:distributed-consistency} shows that
both theorems hold for distributed histories.\cxbar{86}{0.4 0.4 1}{e}

\color{black}\cxtag{87}{N/S}The same condition governs legality.

\color{black}\begin{proposition}[Output-legality boundary]
\label{prop:exposure-boundary}
\cxtag{88}{N/M}\cxbar{88}{0.4 0.4 1}{s}An \spectype satisfies Lifted Monotonicity if and only
if every output at every state reachable without \commitment is legal.\cxbar{88}{0.4 0.4 1}{e}
\end{proposition}

\color{orange!90!black}\begin{proof}[Proof sketch]
\cxtag{89}{O/L}\cxbar{89}{0.9 0.45 0}{s}Under the common-extension condition, consistency supplies the common \outnoun
the legal-output test asks for.  Conversely, the legal-output test at canonical
states, one row at a time, builds the common \outnoun the condition
asks for.  \fullproof\cxbar{89}{0.9 0.45 0}{e}
\end{proof}

\color{black}\cxtag{90}{N/S}For functions, both are monotonicity, so Complete CALM and
Proposition~\ref{prop:exposure-boundary} reduce to Functional CALM and
Proposition~\ref{prop:functional-exposure-boundary}.

% ======================================================================
\section{Distributed Execution}
\label{sec:distributed-execution}

\color{black}\cxtag{91}{N/S}This section gives Functional CALM and Complete CALM a distributed reading,
including a convergence guarantee for functional specs.  We then define
systems and relate their runs to canonical runs; Complete CAP follows.

\subsection{Distributed histories and consistent cuts}
\label{sec:histories}
\label{sec:distributed-consistency}

\color{blue!75!black}\cxtag{92}{B/M}\cxbar{92}{0.4 0.4 1}{s}Until now, runs have been independent: each state changes only through that
run's input and output events.  We now place one such run at each participant in
a finite set $P$.  Participant $p$ starts in the canonical state
$\sigma_P(i_p)$ and still changes its state only through its own specification
events.  The new ingredient is communication.  Besides learning input directly
from the environment, a participant can learn cumulative input from another
participant.\cxbar{92}{0.4 0.4 1}{e}

\color{blue!75!black}\cxtag{93}{B/L}\cxbar{93}{0.9 0.45 0}{s}Participants exchange messages by $\mathsf{send}(m)$ and
$\mathsf{receive}(m)$.  A \emph{message} $m=(j,d)$ carries a cumulative input
$j$, or $\varnothing$ if it carries none, and arbitrary application-specific
data $d$. For $d$, one can envision timestamps, version vectors, state identities, etc.  A receive does not change the specification state; a participant at
$(i,R)$ that receives $(j,d)$ may incorporate $j$ by its own
$\mathsf{input}(k)$ for a legal common \indesc $k$ of $i$ and $j$.  Data can
influence a system's subsequent actions (\S\ref{sec:systems}), but never
directly changes the state.\cxbar{93}{0.9 0.45 0}{e}

\color{black}\cxtag{94}{N/M}\cxbar{94}{0.4 0.4 1}{s}Send and receive events connect the participants' locally ordered runs.  Adding
an edge from each send to its matching receive gives the standard model of a
distributed history; the transitive closure of these edges and the local orders
is happens-before.\cxbar{94}{0.4 0.4 1}{e}

\color{orange!90!black}\begin{definition}[Distributed history, history event, happens-before, consistent cut]
\label{def:distributed-history}
\cxtag{95}{O/XL}\cxbar{95}{0.8 0 0}{s}A \emph{distributed history} $H$ is a finite set of
\emph{history events}, each an occurrence at one participant of a
specification event, $\mathsf{send}(m)$, or $\mathsf{receive}(m)$, with
a total local order on each participant's events, a matching send at another
participant for each receive, and, for each participant $p$, specification
events that in local order form a defined transition sequence from
$\sigma_P(i_p)$.  A \emph{consistent cut} $C$ is a
subset of history events closed under
happens-before~\cite{lamport1978time,chandy1985distributed} (a term
unrelated to our consistency).\cxbar{95}{0.8 0 0}{e}
\end{definition}

\color{blue!75!black}\cxtag{96}{B/M}\cxbar{96}{0.4 0.4 1}{s}At a consistent cut $C$, the specification events of $C$ at participant $p$,
in local order, form its \emph{participant run}, ending in some state
$\sigma_p$.  A \emph{continuation} at $C$ extends the history events of $C$
with further history events.  For a common \indesc $j$ of the participants'
current inputs, confluence asks whether some continuation brings every
participant to the same state $(j,R)$.\cxbar{96}{0.4 0.4 1}{e}

\color{black}\begin{corollary}[Complete CALM at consistent cuts]
\label{cor:distributed-calm}
\cxtag{97}{N/M}\cxbar{97}{0.4 0.4 1}{s}A spec satisfies Lifted Monotonicity if and only if, in every finite
distributed history, every consistent cut is confluent at every common
\indesc of the participants' current inputs.\cxbar{97}{0.4 0.4 1}{e}
\end{corollary}

\color{black}\begin{proof}[Proof sketch]
\cxtag{98}{N/M}\cxbar{98}{0.4 0.4 1}{s}The participant runs at a cut are a finite family of runs, so Complete CALM
gives the forward direction.  Conversely, any finite family of runs can be
placed at participants that exchange no messages.  \fullproof\cxbar{98}{0.4 0.4 1}{e}
\end{proof}

\color{red!80!black}\cxtag{99}{R/L}\cxbar{99}{0.9 0.45 0}{s}The converse uses unrestricted placement; the placement-sensitive variant is
formalized in \extended.\cxbar{99}{0.9 0.45 0}{e}

\color{blue!75!black}\cxtag{100}{B/L}\cxbar{100}{0.9 0.45 0}{s}The forward direction has a concrete witness.  For participant runs with common
\indesc $j$, the survivor test supplies a $u\in\fiber{j}$ that survives in every
endpoint; by Lemma~\ref{lem:run-residual}, $u$ \outrel everything they have
emitted.  At each participant, $\mathsf{input}(j)$ followed by
$\mathsf{output}(u)$ reaches the same state without \commitment
(Proposition~\ref{prop:exposure-boundary}).  This is an existence witness in the
transition system; choosing $u$, communicating it, and scheduling the
transitions belong to the operational model (\S\ref{sec:scope}).\cxbar{100}{0.9 0.45 0}{e}

\color{black}\cxtag{101}{N/M}\cxbar{101}{0.4 0.4 1}{s}For functional specs,
Corollary~\ref{cor:state-by-input} also gives convergence: participants with
the same current input are in the same state, regardless of message paths,
arrival order, or outputs emitted.  This is the convergence that strong eventual consistency~\cite{shapiro2011crdt}
asks of replicated state.  The semilattice merge of CRDTs suffices for it but
is not needed: it follows directly from any consistent functional
spec.\cxbar{101}{0.4 0.4 1}{e}

\subsection{Systems and canonical runs}
\label{sec:systems}

\color{black}\cxtag{102}{N/S}A spec describes which behaviors are permitted, but not how a system chooses
among them.  We now model those choices, and show that a system's inputs and
outputs always remain those of a canonical run.

\begin{figure}[t]
  \centering
  \small
  \setlength{\tabcolsep}{5pt}
  \begin{tabular}[t]{@{}l l@{}}
    \multicolumn{2}{@{}l}{(h) 2PL: $A$ locks ahead of $B$}\\
    \toprule
    $\mathrm{in}$ & $\mathrm{out}$ \\ \midrule
    $\{A,B,C\}$ & \sh{$[A,B,C]$} \\
    $\{A,B,C\}$ & \sh{$[A,C,B]$} \\
    $\{A,B,C\}$ & \sh{$[C,A,B]$} \\
    \grayrow{$\{A,B,C\}$} & \grayrow{$[B,A,C]$} \\ \strikerow
    \grayrow{$\{A,B,C\}$} & \grayrow{$[B,C,A]$} \\ \strikerow
    \grayrow{$\{A,B,C\}$} & \grayrow{$[C,B,A]$} \\ \strikerow
    \bottomrule
  \end{tabular}
  \hfill
  \begin{tabular}[t]{@{}l l@{}}
    \multicolumn{2}{@{}l}{(i) MVTO: $B$ receives the first timestamp}\\
    \toprule
    $\mathrm{in}$ & $\mathrm{out}$ \\ \midrule
    $\{A,B,C\}$ & \sh{$[B,A,C]$} \\
    $\{A,B,C\}$ & \sh{$[B,C,A]$} \\
    \grayrow{$\{A,B,C\}$} & \grayrow{$[A,B,C]$} \\ \strikerow
    \grayrow{$\{A,B,C\}$} & \grayrow{$[A,C,B]$} \\ \strikerow
    \grayrow{$\{A,B,C\}$} & \grayrow{$[C,A,B]$} \\ \strikerow
    \grayrow{$\{A,B,C\}$} & \grayrow{$[C,B,A]$} \\ \strikerow
    \bottomrule
  \end{tabular}
  \caption{Serial orders of three transactions after one \commitment:
  (h) 2PL orders $A$ before $B$; (i) MVTO places $B$ first.  Only complete rows without aborts or reads
  are shown.}
  \label{fig:cc-commitments}
\end{figure}

\color{blue!75!black}\cxtag{103}{B/L}\cxbar{103}{0.9 0.45 0}{s}Sections~\ref{sec:functional-calm}--\ref{sec:complete-calm} used only the
canonical \commitment $x_o$ that an output requires.  Systems may commit
earlier or remove more: the registrar may close the polls before announcing a
result, and a lock manager may fix an order before emitting it.  We therefore
let systems apply their own \commitments, but only in a third kind of event
alongside $\mathsf{input}(j)$ and $\mathsf{output}(o)$: a \emph{\commitment event}
$\mathsf{apply}(\gamma)$, which applies a sequence $\gamma$ of \commitments
and emits nothing.  A \commitment event may precede an output event, which then
has less to remove.\cxbar{103}{0.9 0.45 0}{e}
% As before, at any input, removing rows rules out their outputs;
% removing every row at that input rules out the input as well.

\color{blue!75!black}\cxtag{104}{B/M}\cxbar{104}{0.4 0.4 1}{s}A \emph{system} consists of an I/O spec, a \emph{\commitment basis} $\mathcal B$ and a
possibly nondeterministic \emph{strategy} per participant that picks its next output,
\commitment, send, or receive event, applying only
\commitments from $\mathcal B$.  The full paper gives these definitions formally.
Coordination mechanisms live in the system: its
strategies describe how participants use them, and its basis records the
\commitments they can enforce.\cxbar{104}{0.4 0.4 1}{e}

\color{black}\cxtag{105}{N/M}\cxbar{105}{0.4 0.4 1}{s}Two familiar concurrency control examples make this concrete. Two-phase locking
(2PL) orders transactions pairwise through lock requests: when $B$ waits for a
lock $A$ holds, $\mathsf{order}(A,B)$ fixes $A$ before $B$ and nothing else.
Multiversion timestamp ordering (MVTO)~\cite{bernstein1981concurrency} orders transactions totally through timestamp
assignment: $\mathsf{stamp}(B)$ fixes $B$ before any transaction not yet
stamped. Both are \commitments, but each leaves different possibilities
afterward (Figure~\ref{fig:cc-commitments}).\cxbar{105}{0.4 0.4 1}{e}

\color{black}\cxtag{106}{N/S}We call the runs of
Sections~\ref{sec:functional-calm}--\ref{sec:complete-calm} \emph{canonical
runs}.  The next lemma relates every system run to one of them.

\color{blue!75!black}\begin{lemma}[Observable Equivalence]
\label{lem:canonical-equivalent}
\cxtag{107}{B/L}\cxbar{107}{0.9 0.45 0}{s}Dropping the \commitment events from a run under any \commitment basis
yields \emph{its canonical run}, which is \emph{observably equivalent} to it: it accepts
the same inputs and emits the same outputs, and after each event it keeps at
least the original's rows (\extended).\cxbar{107}{0.9 0.45 0}{e}
\end{lemma}

\color{blue!75!black}\cxtag{108}{B/M}\cxbar{108}{0.4 0.4 1}{s}Hence the commitments a spec's outputs require do not depend on their timing or
representation in a system's \commitment basis.\cxbar{108}{0.4 0.4 1}{e}

\color{orange!90!black}\cxtag{109}{O/L}\cxbar{109}{0.9 0.45 0}{s}Observable equivalence is one-way only: all system runs have observably
equivalent canonical runs, but not necessarily vice versa. Note also that
observable equivalence is a guarantee only on inputs and outputs, not on
remainder tables.  A system run's \commitments may be more aggressive, removing
rows that its canonical run keeps. As a result, a spec can be confluent under
canonical runs, but a system with that spec may not be confluent, e.g., if it
removes a necessary common extension.  The registrar is an example.  After
closing the polls at $ab\cdot$, it cannot accept the late ballot $aba$, although
its canonical run, which only emits $\bot$, can. The CALM and CAP criteria are
about specs and their canonical runs. A system with no \commitment events
inherits their confluence directly, since its runs are canonical runs. For a
system with \commitments of its own, showing that it inherits that confluence
requires an argument about that system, which is outside the scope of this
paper.\cxbar{109}{0.9 0.45 0}{e}

\subsection{The Complete CAP Theorem}
\label{sec:complete-cap}

\color{black}\cxtag{110}{N/S}Brewer conjectured that a distributed system cannot be consistent, available,
and partition tolerant at once~\cite{brewer2000towards}; Gilbert and Lynch
proved it as an impossibility result for a linearizable register~\cite{gilbert2002cap}.
Complete CAP addresses the conjecture more broadly:
which specifications are subject to the CAP
conjecture, and which are not?

\color{blue!75!black}\cxtag{111}{B/M}\cxbar{111}{0.4 0.4 1}{s}To pose the question properly, note that the outputs that each participant emits---and
needs to keep available---are not enough to judge consistency. The spec
must also model what an observer (e.g., a client) can eventually conclude by collecting
outputs from multiple participants.\cxbar{111}{0.4 0.4 1}{e}

\color{orange!90!black}\cxtag{112}{O/L}\cxbar{112}{0.9 0.45 0}{s}Like Gilbert and Lynch, we can use a replicated register as an example.
A register is a simplified key-value store that holds only the
value of a single key, which suffices for reasoning about properties of a single object.
Consider a racing read and write that go to different
participants. The register is initialized to $v_0$. Participant $p$ receives
write $w(v)$ and emits its acknowledgment; concurrently, participant $q$
receives the read $r$ and returns $v_0$.  A client may later collect both the
write acknowledgment and read answer $v_0$.  A linearizable register spec has a
row that allows this situation: that pair of requests forms an input, and the
combined responses are an output. That row reflects that a linearization may put the
concurrent read before the write.  However, no participant emits this combined
output, nor is the combined output needed for availability. Nonetheless, it \outrel
both local outputs, provides a proof of their consistency, and must be in the spec.  Had the write finished
before the read began, no such combined output would exist; we return to this
case at the end of the section.\cxbar{112}{0.9 0.45 0}{e}

\color{blue!75!black}\cxtag{113}{B/M}\cxbar{113}{0.4 0.4 1}{s}We therefore pair a spec with an \emph{explicit possibility relation}
$E\subseteq P$, or explicit relation for short: the rows whose outputs
participants actually emit.  Outputs in
$P\setminus E$ are not emitted by participants;
they remain permitted in $P$ to capture implicit combinations of emitted outputs---e.g.,
as could be constructed by external observers or clients making sense of explicit outputs
across participants.\cxbar{113}{0.4 0.4 1}{e}

\color{black}\cxtag{114}{N/S}Both availability and consistency are judged on histories in which a partition
separates the participants.

\color{black}\begin{definition}[Partitioned history]
\label{def:partitioned-history}
\cxtag{115}{N/M}\cxbar{115}{0.4 0.4 1}{s}A \emph{partitioned history} divides the participants into two or more
sides and contains no receive whose matching send lies on another
side; participants still receive environment input and emit outputs.\cxbar{115}{0.4 0.4 1}{e}
\end{definition}

\color{black}\begin{definition}[Explicit possibility relation]
\label{def:explicit-relation}
\cxtag{116}{N/S}An \emph{explicit possibility relation} $E\subseteq P$ is a set of rows whose
outputs participants actually emit.  Write $E(i)$ for its outputs at input
$i$.
\end{definition}

\color{blue!75!black}\cxtag{117}{B/M}\cxbar{117}{0.4 0.4 1}{s}For example, consider two-phase commit (2PC).  We focus on the spec for the
voters; the coordinator is an observer of that spec.  We're interested in the
fiber at the input \textsf{prepare(tx)}: how each voter responds to that
prepare message.  Each possible explicit response is a row of $E$.  The overall
\emph{set} of responses the coordinator might observe across voters is in
$P\setminus E$.  The spec includes these as implicit outputs, representing the
possible combined outcomes that respect each voter's response.\cxbar{117}{0.4 0.4 1}{e}

\color{black}\cxtag{118}{N/S}We consider a system, with a strategy for each participant that depends only on its
local history, and an environment that may send any legal input to any
participant.

\color{blue!75!black}\begin{definition}[Availability]
\label{def:partition-notions}
\cxtag{119}{B/L}\cxbar{119}{0.9 0.45 0}{s}A \emph{same-side} continuation of a participant $p$ is one the system can
produce using only events on the side of the partition containing $p$, with no
new input and no message from another side.  The system has
\emph{Availability for $E$} when, in every partitioned history it produces,
each participant $p$ in state $(i,R)$ has, for each output $o\in R(i)\cap E(i)$, some same-side continuation in which $p$ emits $o$.
\emph{Availability} is shorthand for Availability
for $P$.\cxbar{119}{0.9 0.45 0}{e}
\end{definition}

\color{black}\begin{definition}[Partition Consistency]
\label{def:partition-consistency}
\cxtag{120}{N/M}\cxbar{120}{0.4 0.4 1}{s}The system has \emph{Partition Consistency} when the final cut of every
partitioned history it produces is confluent at every common \indesc of the
participants' current inputs.\cxbar{120}{0.4 0.4 1}{e}
\end{definition}

\color{black}\cxtag{121}{N/S}The \emph{canonical system for $E$} is the system whose participants follow
canonical runs and emit only outputs of $E$.

\color{black}\begin{definition}[Escaping CAP]
\label{def:escapes-cap}
\cxtag{122}{N/M}\cxbar{122}{0.4 0.4 1}{s}A spec \emph{escapes CAP} for $E$ when its canonical system for $E$ has
Availability for $E$ and Partition Consistency.\cxbar{122}{0.4 0.4 1}{e}
\end{definition}

\color{black}\cxtag{123}{N/M}\cxbar{123}{0.4 0.4 1}{s}The canonical system's participants never wait: each can emit any output of
$E$ that its state allows, so it always has Availability for $E$.  The
question is whether their answers stay consistent across sides.  Complete CAP
says they do exactly when outputs of $E$ emitted on different sides have a
common \outnoun in the spec at every common \indesc of their inputs.\cxbar{123}{0.4 0.4 1}{e}

\color{black}\begin{definition}[Lifted Monotonicity for $E$]
\label{def:lm-for-A}
\cxtag{124}{N/M}\cxbar{124}{0.4 0.4 1}{s}$\Spec$ satisfies \emph{Lifted Monotonicity for $E$} when every finite set of rows of $E$ with
a common \indesc $j$ has a common \outnoun in $\fiber{j}$.\cxbar{124}{0.4 0.4 1}{e}
\end{definition}

\color{black}\begin{theorem}[Complete CAP]
\label{thm:complete-cap}
\cxtag{125}{N/S}A spec escapes CAP for $E$ if and only if it satisfies Lifted Monotonicity for
$E$.
\end{theorem}

\color{blue!75!black}\cxtag{126}{B/S}When $E=P$, we can say simply that a spec escapes CAP exactly when it
satisfies Lifted Monotonicity.

\color{blue!75!black}\begin{proof}[Proof sketch]
\cxtag{127}{B/L}\cxbar{127}{0.9 0.45 0}{s}If Lifted Monotonicity for $E$ holds, the outputs emitted on all sides have a
common \outnoun $u$ at each common \indesc $j$, and every participant still has
the row $(j,u)$, so the participants can meet there.  Conversely, put one
participant per row of $E$ on its own side.  Each emits its row's output, and
Partition Consistency gives those outputs a common \outnoun.  \fullproof\cxbar{127}{0.9 0.45 0}{e}
\end{proof}

\color{black}\begin{corollary}
\label{cor:complete-cap-any-system}
\cxtag{128}{N/M}\cxbar{128}{0.4 0.4 1}{s}No system does better than its spec: if some system has Availability for $E$
and Partition Consistency, its spec escapes CAP for $E$.\cxbar{128}{0.4 0.4 1}{e}
\end{corollary}

\color{blue!75!black}\begin{proof}[Proof sketch]
\cxtag{129}{B/L}\cxbar{129}{0.9 0.45 0}{s}If Lifted Monotonicity for $E$ fails, some outputs of $E$ have no common
\outnoun at a common \indesc $j$.  A system with Availability for $E$ can emit
them on separate sides.  By Observable Equivalence, each participant keeps no
more rows than its canonical run, so no participant has a row at $j$ that
\outrel all of them.\cxbar{129}{0.9 0.45 0}{e}
\end{proof}

\color{blue!75!black}\cxtag{130}{B/S}For $E=P$, the canonical system needs no \commitment.

\color{blue!75!black}\cxtag{131}{B/M}\cxbar{131}{0.4 0.4 1}{s}Return to the voters' spec.  Each voter decides whether it can commit and emits
its vote, yes or no.  That decision is a \commitment: a voter that votes yes
rules out every set of votes in which it votes no.  But it is a single-voter
\commitment involving no communication, so the voter can make it and produce
output during a partition.  Each vote is cast by a different voter, so votes
emitted on different sides of a partition never conflict (the full paper formalizes this placement as a deployment~\extended), and every
set of votes an observer may collect is consistent.\cxbar{131}{0.4 0.4 1}{e}

\color{black}\cxtag{132}{N/M}\cxbar{132}{0.4 0.4 1}{s}A partition does not change the intrinsic \commitments a spec needs; rather,
it prevents separated participants from jointly enforcing those \commitments
via coordination.  When Lifted Monotonicity
for $E$ fails, some output of $E$ waits for another side, sacrificing
Availability, or is emitted anyway, sacrificing Partition Consistency.\cxbar{132}{0.4 0.4 1}{e}

\paragraph{Gilbert and Lynch.}
\color{blue!75!black}\cxtag{133}{B/L}\cxbar{133}{0.9 0.45 0}{s}Their theorem is a case of Complete CAP.  We return to the register example, but
add the requirement of linearizability regarding real-time precedence: if one
operation finishes before another begins, every linearization must put the
first before the second~\cite{herlihy1990linearizability}.  Now let the write
finish before the read begins.  Every linearization must then put the write
first, so the read must return $v$, and no output \outrel $q$'s answer $v_0$.
Since Availability requires $p$ to acknowledge the write and $q$ to answer
$v_0$, the spec does not satisfy Lifted Monotonicity for those outputs.  By
Corollary~\ref{cor:complete-cap-any-system}, no system can keep both outputs
available and stay consistent across a partition.  This is Gilbert and Lynch's
theorem.
The full paper takes precedence from real time, extends
Complete CAP to systems that choose their responses when responses are not
forced, and derives
the theorem in Gilbert and Lynch's asynchronous model~\cite{hellerstein2026completecalmfull}.\cxbar{133}{0.9 0.45 0}{e}

% ======================================================================
\section{Composition and Proper Coordination}
\label{sec:composition}

\color{black}\cxtag{134}{N/M}\cxbar{134}{0.4 0.4 1}{s}A replicated key-value store poses a puzzle.  With mutable values, it is
non-monotone and inconsistent~\cite{alvaro2011bloom}. The pragmatic solution is
to impose a consensus-driven log upstream of the replicas, \`a la
Multi-Paxos~\cite{lamport2001paxos}, so all replicas execute the same commands
in the same order. But consensus is itself non-monotone: when two commands
arrive concurrently, either may be appended next, and the two alternative orders
have no common \outnoun.  In essence, the standard solution to non-monotonicity
is more non-monotonicity! Why does it work?\cxbar{134}{0.4 0.4 1}{e}

\color{black}\cxtag{135}{N/M}\cxbar{135}{0.4 0.4 1}{s}That puzzle is explained by our framework.  As we'll see, two specs compose by equijoin,
and the composition of two monotone specs is also monotone. Moreover, for non-monotone specs,
as in our puzzle, we
can reason about which component actually makes \commitments under composition.  When the
upstream component makes every \commitment the composition needs, and the
downstream needs none on the inputs the upstream offers, we say the pipeline's
coordination is \emph{proper}
(\S\ref{sec:proper-coordination}).\cxbar{135}{0.4 0.4 1}{e}
\label{sec:sequential-composition}

\color{blue!75!black}\cxtag{136}{B/L}\cxbar{136}{0.9 0.45 0}{s}Formally, consider two \spectype[s]
${\Spec_1=(I,\InputOrd,M,\preceq_M,P_1)}$ and
${\Spec_2=(M,\preceq_M,O,\OutputOrd,P_2)}$
sharing the intermediate ordered domain $(M,\preceq_M)$.  Their \emph{sequential
composition} has table
\begin{equation}
\begin{split}
  P_2\circ P_1
  &=\pi_{I,O}(P_1\Join_M P_2)\\
  &=\{(i,o):\exists m.\ (i,m)\in P_1\land(m,o)\in P_2\}
\end{split}
  \label{eq:spec-composition}
\end{equation}
where $\pi$ and $\Join$ are relational projection and equijoin.  Since $P_1$
is total on $I$ and $P_2$ is total on $M$, every $i\in I$ has a nonempty fiber
in $P_2\circ P_1$.  When the two stages
run as a pipeline, $\Spec_2$ takes the outputs of $\Spec_1$ as its inputs, and
we say that $\Spec_1$ \emph{guards} $\Spec_2$.\cxbar{136}{0.9 0.45 0}{e}

\color{black}\begin{lemma}[Composition]
\label{lem:composition}
\cxtag{137}{N/S}If $\Spec_1$ and $\Spec_2$ satisfy Lifted Monotonicity, then so does the
\spectype with table $P_2\circ P_1$.
\end{lemma}

\color{orange!90!black}\begin{proof}[Proof sketch]
\cxtag{138}{O/L}\cxbar{138}{0.9 0.45 0}{s}Take rows of $P_2\circ P_1$ whose inputs have a common \indesc $j$.  The
common-extension condition for $\Spec_1$ gives an intermediate value at $j$
that \outrel each row's, and that value is a common \indesc
at which the condition for $\Spec_2$ supplies the common \outnoun.  \fullproof\cxbar{138}{0.9 0.45 0}{e}
\end{proof}

\color{black}\cxtag{139}{N/S}By associativity the lemma extends to finite pipelines.

\subsection{Proper Coordination}
\label{sec:proper-coordination}

\color{black}\cxtag{140}{N/M}\cxbar{140}{0.4 0.4 1}{s}Return to the replicated store of the opening.  The key-value store is not
monotone on its own, but it is monotone on what the log offers: a log fixes
the order of commands, and a longer log only extends the responses.  The log
issues a \commitment to offer only such inputs, freeing the store from making any
\commitments.\cxbar{140}{0.4 0.4 1}{e}

\color{blue!75!black}\cxtag{141}{B/M}\cxbar{141}{0.4 0.4 1}{s}More generally, let $\Spec_1$, which guards $\Spec_2$, be run by a system with any basis, and
let $\Spec_2$ run under a system of its own.  Write $\Spec_2|_{\Spec_1}$ for $\Spec_2$ restricted to the outputs of
$\Spec_1$.\cxbar{141}{0.4 0.4 1}{e}

\color{black}\begin{definition}[Proper Coordination]
\label{def:proper-coordination}
\cxtag{142}{N/M}\cxbar{142}{0.4 0.4 1}{s}A pipeline in which $\Spec_1$ guards $\Spec_2$ is \emph{proper} when
$\Spec_2|_{\Spec_1}$ needs no \commitment: every output at every state it
reaches without \commitment is legal.\cxbar{142}{0.4 0.4 1}{e}
\end{definition}

\color{blue!75!black}\begin{theorem}[Proper Coordination]
\label{thm:proper-coordination}
\cxtag{143}{B/L}\cxbar{143}{0.9 0.45 0}{s}Fix any system for $\Spec_1$.
\begin{enumerate}
  \item The pipeline is proper exactly when $\Spec_2|_{\Spec_1}$ satisfies
        Lifted Monotonicity.
  \item If it is proper, $\Spec_2|_{\Spec_1}$ can take each output of a run of
  $\Spec_1$ as its next input and answer it without \commitment. The result is a
  run of $P_2\circ P_1$, and composite runs are confluent whenever the
  underlying runs of $\Spec_1$ are.
  \item If it is proper and $P_2\circ P_1$ violates Lifted Monotonicity,
        $\Spec_1$ cannot emit some output without a \commitment.\cxbar{143}{0.9 0.45 0}{e}
\end{enumerate}
\end{theorem}

\color{black}\begin{proof}[Proof sketch] \cxtag{144}{N/M}\cxbar{144}{0.4 0.4 1}{s}(1) is Complete CALM and the output-legality
boundary (Theorem~\ref{thm:complete-calm},
Proposition~\ref{prop:exposure-boundary}) applied to $\Spec_2|_{\Spec_1}$. The
full proof of (1) and the proofs of (2) and (3) are in \extended.\cxbar{144}{0.4 0.4 1}{e}
\end{proof}

\color{blue!75!black}\cxtag{145}{B/M}\cxbar{145}{0.4 0.4 1}{s}The theorem holds whatever basis $\Spec_1$ uses, and a proper pipeline is
itself a guard for a further stage.  Locks in
the key-value store would be pure overhead: by (1), the store restricted to
logs needs no \commitment.\cxbar{145}{0.4 0.4 1}{e}

\subsection{Composite CALM}
\label{sec:composite-calm}

\color{orange!90!black}\cxtag{146}{O/L}\cxbar{146}{0.9 0.45 0}{s}Proper coordination answers our opening puzzle by capturing how a pipeline's
\commitment obligations can be discharged by an upstream component.  We now turn
to a simpler puzzle. Consider the spec for an incrementing counter that is
capped at $9$. Now compose it with a downstream spec for a double-digit
threshold test: that spec outputs \textsf{false} for values up through $9$, and
\textsf{true} as the input grows further. In the
literature~\cite{kuper2013lvars}, the output order is the usual boolean lattice,
so \textsf{false} is a lower bound; only \textsf{true} is a guarantee.
Suppose we want the downstream spec to always give a final answer; then the
output order is equality, and the test cannot emit \textsf{false} without making
a \commitment. Nonetheless, when guarded by the capped counter the composition becomes
\commitment-free. Can we capture that in our framework? When do composed specs
erase the need for \commitment entirely?\cxbar{146}{0.9 0.45 0}{e}

\color{orange!90!black}\cxtag{147}{O/L}\cxbar{147}{0.9 0.45 0}{s}By Complete CALM, $P_2\circ P_1$ is consistent exactly when it
satisfies Lifted Monotonicity, and then no output needs a \commitment.
A composite can be jointly monotone, with no stage needing a
\commitment, in two ways.  First, if every stage is monotone,
as when a monotone client query observes CRDT state
(\S\ref{sec:crdt-instance}), then Lemma~\ref{lem:composition} applies.
Second, the upstream can filter out the inputs on which the downstream is not
monotone: the counter capped at $9$ only produces downstream inputs that match
with \textsf{false} as an output, so the guarded downstream spec is consistent---under any output order,
including equality.
This is proper coordination in which the upstream itself
needs no \commitment.
There are additional scenarios resulting in monotone composition beyond these cases
(see \extended).\cxbar{147}{0.9 0.45 0}{e}

\color{black}\cxtag{148}{N/M}\cxbar{148}{0.4 0.4 1}{s}The full paper extends these results to more complex
compositions than sequential stages, including
stages with inputs of their own and arbitrary acyclic graphs of stages, checked node
by node much as types are~\cite{hellerstein2026completecalmfull}.\cxbar{148}{0.4 0.4 1}{e}

% ======================================================================
\section{Related Work and Mappings}
\label{sec:instances}
\label{sec:related}

\color{blue!75!black}\cxtag{149}{B/M}\cxbar{149}{0.4 0.4 1}{s}Complete CALM provides a unified lens for many results across literatures.
For the key related ideas, we provide mappings here to our specification model.
Each setting declares an input, an output, and their orders before
applying either CALM theorem.  The full paper~\cite{hellerstein2026completecalmfull}
gives each mapping in detail, with deployments, response
obligations, and \commitment bases for the transaction protocols 2PL and MVTO,
and applications to distributed-systems protocols such as registers, $k$-set agreement, renaming, and snapshots.\cxbar{149}{0.4 0.4 1}{e}

\subsection{Query monotonicity}
\label{sec:query-instance}

\color{orange!90!black}\cxtag{150}{O/XL}\cxbar{150}{0.8 0 0}{s}As previously discussed (\S\ref{sec:functional-calm}),
the CALM Theorem of Ameloot et al.\ proves an operational result for
relational languages on transducer networks: coordination-free evaluation is possible
iff a program is monotone,
under its data-placement and execution assumptions~\cite{ameloot2013relational,ameloot2015weaker}.
Functional CALM does not replace the details of that theorem.  It abstracts
the result to a property that can be compared with specs outside
Datalog. Subsequent extensions to the original CALM Theorem become changes to the input domain
and its order.
When transducers know
how data is placed, more queries become coordination-free, among them the
non-monotone win-move query~\cite{zinn2012winmove}, and Ameloot et al.\
characterized the new classes by domain-distinct and domain-disjoint
monotonicity~\cite{ameloot2015weaker,zinn2012weak}.  Each is Functional CALM
under an input order in which $I\cup J$ \inrel $I$ only for the $J$ that the
placement policy leaves a node unable to rule out.  Functional CALM takes the
order as given.\cxbar{150}{0.8 0 0}{e}

\color{blue!75!black}\cxtag{151}{B/L}\cxbar{151}{0.9 0.45 0}{s}Baccaert and Ketsman's \emph{behaviors} map located input facts to sets of
permitted located outputs~\cite{baccaert2023dcbq,baccaert2026spectrum}; under
inclusion, a behavior is an \spectype.  Their characterization,
\emph{node monotonicity}, is the common-extension condition for two rows at
disjoint networks, and Complete CAP extends it to overlapping networks
(see \extended).\cxbar{151}{0.9 0.45 0}{e}

\subsection{CRDTs: monotone queries over convergent state}
\label{sec:crdt-instance}

\color{blue!75!black}\cxtag{152}{B/L}\cxbar{152}{0.9 0.45 0}{s}A CRDT guarantees that replicas receiving the same updates converge to
equivalent states, but the Strong Eventual Consistency property it offers
provides no guarantee on what the environment may observe.  Laddad et al.\
proposed monotonicity as the right test for safe observation from CRDT
state~\cite{laddad2022keep}.  We make the proposed test precise with
Functional CALM.  Let $(D,\sqsubseteq,\sqcup)$ be the state join-semilattice.
The input is the finite set $U$ of immutable payloads received so far, ordered
by inclusion, and the replica's state is $s(U)=s_0\sqcup\bigsqcup U$.  A
client observation $q:D\to(Y,\preceq_Y)$ gives the client-visible
spec $q\circ s$.  By Functional CALM, $q\circ s$ is consistent
exactly when $U\subseteq V$ implies $q(s(U))\preceq_Y q(s(V))$.  Monotonicity
of $q$ on all of $D$ suffices, but $q$ need only be monotone between reachable
states.  The query $\mathit{Add}\setminus\mathit{Del}$ over a two-phase set
fails the test: a deleted item grows $U$ and shrinks the answer.  CRDT states
form a join-semilattice, but $q$ needs less: Complete CALM requires
only \emph{some} common \outnoun of query outputs, not a least one.\cxbar{152}{0.9 0.45 0}{e}

\subsection{Highly available transactions}
\label{sec:hat-instance}

\color{black}\cxtag{153}{N/M}\cxbar{153}{0.4 0.4 1}{s}We take the assumptions of Bailis et al.\ as given~\cite{bailis2013hat}:
items are partially replicated; a client sends each operation to some replica
of its item; a transaction's writes carry a unique client timestamp, last
writer wins; the system never aborts on its own; and a replica answers a read
with the highest-timestamped committed write of the item it has received.\cxbar{153}{0.4 0.4 1}{e}

\color{blue!75!black}\begin{definition}[HAT specification]
\label{def:hat-spec}
\cxtag{154}{B/XL}\cxbar{154}{0.8 0 0}{s}An input is a history ordered by happens-before: reads tagged with their
replica, commits with their timestamp and write set, and deliveries of
committed writes to replicas.  Inputs are ordered by
downward-closed extension, and outputs are sets of read answers, ordered by
inclusion.  For an isolation level $L$, defined by the phenomena it
proscribes~\cite{adya1999weakconsistency}, let $\Spec_L$ have table $P_L$,
which permits an output when it
extends to a history that answers every read, orders versions by timestamp,
and exhibits none of them.  The explicit relation $E$ holds, at each
input, the answers the read rule gives.\cxbar{154}{0.8 0 0}{e}
\end{definition}

\color{black}\begin{theorem}[HAT levels]
\label{thm:hat-levels}
\cxtag{155}{N/L}\cxbar{155}{0.9 0.45 0}{s}Suppose each transaction's timestamp exceeds that of every version it reads.
For read uncommitted and read committed, $E\subseteq P_L$ and $\Spec_L$
satisfies Lifted Monotonicity for $E$, so by Complete CAP these levels stay
available and consistent through partitions.  For any level that proscribes
Lost Update, including snapshot isolation, repeatable read, and
serializability, no explicit relation that answers every read satisfies
Lifted Monotonicity for it.\cxbar{155}{0.9 0.45 0}{e}
\end{theorem}
\proofinfull

\color{orange!90!black}\cxtag{156}{O/L}\cxbar{156}{0.9 0.45 0}{s}The condition is needed: with timestamps fixed at transaction
begin, as in the pseudocode of~\cite{bailis2013hatextended}, a transaction
that reads a higher-timestamped version and overwrites it forms a G1c
cycle.\cxbar{156}{0.9 0.45 0}{e}

\color{black}\cxtag{157}{N/M}\cxbar{157}{0.4 0.4 1}{s}Mahajan et al.\ showed that no consistency semantics stronger than real-time
causal consistency can be enforced in an always-available, one-way convergent
system~\cite{mahajan2011consistency}.  With real-time precedence recorded in
the input, both halves of their result follow from Complete CAP (\extended).\cxbar{157}{0.4 0.4 1}{e}

\subsection{Invariant confluence}
\label{sec:iconfluence-instance}

\color{blue!75!black}\cxtag{158}{B/XL}\cxbar{158}{0.8 0 0}{s}In the model of Bailis et al.~\cite{bailis2014coordination}, a database state
is a set of versions, replicas merge states by union, and a set $T$ of
transactions is \emph{invariant confluent} for an invariant $\mathcal I$ when
the merge of any two reachable states satisfies $\mathcal I$.  We construct a
spec $\Spec_{\mathcal I}$ whose inputs are locally valid
\emph{transaction histories}, partial orders of transactions in which one
precedes another when the second's replica saw the first, ordered
by extension.  Outputs are database states, ordered by inclusion;
$P_{\mathcal I}$ permits at $h$ the state of any downward-closed
sub-history of $h$ whose merges all satisfy $\mathcal I$.\cxbar{158}{0.8 0 0}{e}

\color{black}\begin{proposition}[Invariant confluence and Lifted Monotonicity]
\label{prop:iconfluence-lm}
\cxtag{159}{N/S}$T$ is invariant confluent if and only if $\Spec_{\mathcal I}$ satisfies Lifted
Monotonicity.
\end{proposition}
\proofinfull

\color{orange!90!black}\cxtag{160}{O/L}\cxbar{160}{0.9 0.45 0}{s}In Bailis et al.'s payroll example, two
replicas each assign ID~7 to a new employee; the history in which neither
creation saw the other \inrel both one-creation histories, but no state it
permits \outrel both.\cxbar{160}{0.9 0.45 0}{e}

\subsection{Operational models and scope}
\label{sec:scope}

\color{blue!75!black}\cxtag{161}{B/M}\cxbar{161}{0.4 0.4 1}{s}Automata are a traditional way to model operational behavior
abstractly~\cite{lynch1987io,lynch1996distributed}.  Our model keeps only the
inputs and outputs of that behavior and drops the operational details.  The
mapping is straightforward: for each input, record the outputs of the runs that
consume it.  Two standard approaches then apply, both used widely and with
success.  Model checkers explore runs on bounded
inputs~\cite{flanagan2005dpor,newcombe2015aws};
their input/output pairs become a finite spec in our model that can be tested for
violations of Lifted Monotonicity.  The only additional requirement is on the
checker's bounds: any input it reports on must have had all its executions
explored. The more comprehensive approach is to use
formal proofs to declare a spec and show, by the standard obligations
of refinement~\cite{abadi1991refinement,lynch1995forward}, that it represents
the machine: every execution's output is dominated by a declared one, and every
declared row is produced by some execution.  Such proofs are routinely
mechanized~\cite{hawblitzel2015ironfleet}, and Lifted Monotonicity is then
checked on the declared spec; \extended{} gives the details.  Work that
models systems as state machines therefore maps naturally into our model, where
Complete CALM assesses its consistency directly.\cxbar{161}{0.4 0.4 1}{e}

  \color{black}\cxtag{162}{N/M}\cxbar{162}{0.4 0.4 1}{s}Our conditions are semantic.  Confluence guarantees that
a common \outnoun exists, not that a deployed system reaches it; protocols,
liveness, and cost belong to the operational model in which a spec
is realized.  Free termination likewise corresponds to a determined canonical
state in our model, but Power et al.\ additionally study fairness and eventual
response, liveness concerns whose semantic generalization remains future
work~\cite{power2025freetermination}.\cxbar{162}{0.4 0.4 1}{e}

% ======================================================================
\section{Discussion}
\label{sec:discussion}

\color{black}\cxtag{163}{N/M}\cxbar{163}{0.4 0.4 1}{s}The computing literature is rich with coordination mechanisms for managing
concurrency within and across machines. The CALM Conjecture asked a question
still latent at the time: when are these mechanisms required, and why? It drew
on declarative networking~\cite{loo2009declarative}, which, in Codd's tradition
of separating specification from implementation~\cite{codd1970relational}, wrote
the desired effects of protocols as queries to be analyzed and optimized into a
variety of mechanisms. Could a signature property of declarative
languages---monotonicity---be the criterion?\cxbar{163}{0.4 0.4 1}{e}

\color{black}\cxtag{164}{N/M}\cxbar{164}{0.4 0.4 1}{s}The intuition was sharpened by a related friction in the literature on parallel
query processing. The Google MapReduce
paper~\cite{dean2004mapreduce} departed from prior work on parallel query
processing~\cite{wilschut1991dataflow,graefe1990volcano}, conflating data
redistribution (shuffle) with barrier coordination required for a non-monotone
reduce operator. This coupling of unrelated concerns required monotone operators
like join to block in the MapReduce model unnecessarily.  Cost models that count
rounds, from BSP supersteps to MapReduce rounds, carry the same
conflation~\cite{valiant1990bsp,karloff2010mrc}.\cxbar{164}{0.4 0.4 1}{e}

\color{black}\cxtag{165}{N/S}The CALM Conjecture hypothesized that the database community's more judicious
use of barriers could be generalized: coordination---of any kind---is required
if and only if the program is non-monotone. Ameloot et al.\
proved it, but only for relational languages, which to date have
still not become mainstream beyond query processing (more on this below).

\color{blue!75!black}\cxtag{166}{B/M}\cxbar{166}{0.4 0.4 1}{s}This paper moves the declarations one level higher, from queries to views: all
that matters about a spec is the single view mapping a program's inputs
to its output. Eliding the query logic frees the analysis from any particular
language and broadens what can be specified. Prior work focused on implementing
complex programs
declaratively~\cite{loo2005overlays,loo2009declarative,chu2007declarative,alvaro2011bloom}.
This was interesting ``declarative engineering'' with concrete benefits in
analysis and optimization. But the higher-level model here also captures abstract
families of guarantees, like CAP and isolation levels, which have no application or
protocol in mind.\cxbar{166}{0.4 0.4 1}{e}

\color{blue!75!black}\cxtag{167}{B/M}\cxbar{167}{0.4 0.4 1}{s}Lifting the abstraction improves the power of analysis as well. Prior program
analyses became enmeshed in the details of implementations, however
declarative. Bloom~\cite{alvaro2011bloom} flagged syntactic non-monotonicity in
programs and conservatively injected coordination code in front of it.  Yet
coordination code is itself non-monotone, so adding it to a Bloom program only
adds more non-monotone syntax.  Subsequent work moved the question from syntax to declared contracts and
invariants~\cite{sivaramakrishnan2015quelea,kaki2018alone}.  Their tools pick,
for each operation or transaction, the weakest consistency or isolation level
that the store offers and that still satisfies the declared properties.  We
ask what the spec itself requires, independent of any menu of
levels.  By Observable Equivalence (Lemma~\ref{lem:canonical-equivalent}), the answer holds for every mechanism.  Our
composition results need no program logic, and the criteria are exact on
specs, not conservative.\cxbar{167}{0.4 0.4 1}{e}

\color{black}\cxtag{168}{N/M}\cxbar{168}{0.4 0.4 1}{s}A limitation of prior CALM work was that it was constructive on only one side of
its boundary.  Monotone programs need no coordination; for a non-monotone
program, CALM said only that coordination was required somewhere, not what it
had to accomplish or where it belonged.  Complete CALM says more.  A
spec that is not monotone needs specific \commitments: the \commitments
its outputs require.  Every coordination mechanism, from locks to consensus,
is a way of applying these \commitments.  The same analysis says which outputs
need \commitments and whether conflicting choices can arise at another
participant.  That tells an implementation when it needs agreement rather than
a local decision.  Proper coordination
then places the \commitments in a composition: it identifies the stage that
must make them, as the log does in state-machine replication, and certifies
that the stages behind it need none of their own.\cxbar{168}{0.4 0.4 1}{e}

\color{black}\cxtag{169}{N/M}\cxbar{169}{0.4 0.4 1}{s}Underlying these issues is a reassessment of
tradition: what is a declarative specification, and how little does it require
to be useful? The question is timely. Coding agents increasingly write
imperative code, often rife with bugs, which has renewed interest in formal
specification languages like Lean~\cite{demoura2021lean4}. A higher-level
language keeps agents out of the incidental complexity of implementations and
pairs them with proofs of correctness: neurosymbolic
computing~\cite{an2025neurosymbolic}. This is a renewed opportunity for logic,
dataflow, and functional languages. Today these languages typically specify
protocols and programs in detail, as in declarative networking; this paper
points to specification at a higher level still, opening analyses like ours and
potential program synthesis and optimization.\cxbar{169}{0.4 0.4 1}{e}

% ======================================================================
\section{Conclusion}
\label{sec:conclusion}

\color{black}\cxtag{170}{N/M}\cxbar{170}{0.4 0.4 1}{s}A spec abstracts execution details away, and the results here reduce
to a single question: can a participant act on what it knows without being
contradicted by what it learns later?  Mutexes, locks, atomic instructions, and
agreement protocols prevent such contradictions by issuing \commitments.
Functional CALM answers the question for functions: a functional spec
stays confluent as input arrives exactly when it is monotone, and then no
output needs a \commitment.  This generalizes the classic result from logical fact
sets to arbitrary input and output orders.\cxbar{170}{0.4 0.4 1}{e}

\color{orange!90!black}\cxtag{171}{O/L}\cxbar{171}{0.9 0.45 0}{s}Complete CALM extends the answer to specs that offer several outputs
for one input.  The test is Lifted Monotonicity, which combines Relational
Monotonicity with directed fibers.  When it fails, some reachable output
requires a \commitment.  At each input it changes, a \commitment either empties
the fiber or removes particular output rows.  Because \commitments are explicit, the model
can also check where a composition coordinates
(\S\ref{sec:composition}).  A totally ordered log and a
key-value store are both non-monotone at their ordinary specs, yet
when the log guards the store, the log's \commitments are all the composition
needs and the store makes none.  Multi-Paxos is one way to distribute those
\commitments.\cxbar{171}{0.9 0.45 0}{e}

\color{black}\cxtag{172}{N/M}\cxbar{172}{0.4 0.4 1}{s}Complete CAP is the partition form of the same equivalence.  A system can keep
the outputs selected by a spec's explicit relation available, and
consistent, through partitions exactly when they have a common \outnoun at
every common \indesc of their inputs; when the relation selects every output,
this is Lifted Monotonicity.\cxbar{172}{0.4 0.4 1}{e}

\color{black}\cxtag{173}{N/M}\cxbar{173}{0.4 0.4 1}{s}The generality of our specification model opens a range of new questions
regarding the nondeterminism of specs and their \commitments. One direction
extends the boolean criteria of this paper to metrics: a Determination
Complexity theory that assesses \emph{how many} \commitments are intrinsically
needed to determine a spec. Another is Determination Provenance: a
given output---or its absence---is determined not only by inputs, but also by
the \commitments made in a run, which brings formal models of provenance to
concurrent systems. A third direction is to provide general-purpose static
analysis of CALM and proper coordination in a real programming framework.\cxbar{173}{0.4 0.4 1}{e}

\color{black}\cxtag{174}{N/S}\begin{acks}
Thanks to Conor Power, whose persistent questioning of the CALM canon prompted this work.
Much gratitude to Peter Alvaro, Natacha Crooks, Chris Douglas, Tyler Hou, Martin
Kleppmann, Paris Koutris, Edward Lee, Shulu Li, Vincent Liu, David Chu McElroy,
Sarah Morin, and Dan Suciu for helpful feedback. Thanks also to
anonymous reviewers of this and previous submissions for frank feedback that pushed
this work deeper.

\color{black}\cxtag{175}{N/S}Generative AI tools (Infinity, Kiro, Claude and ChatGPT) were used as interactive writing and
reviewing assistants during the preparation of this work, including drafting
and revising text, and checking mathematical arguments with Lean proofs.  All
content was reviewed, validated, and approved by the author.
\end{acks}

\bibliographystyle{ACM-Reference-Format}
\bibliography{references,cidr-extra}

\end{document}
